\documentclass[10pt,titlepage]{article}
\usepackage[letterpaper,margin=0.75in]{geometry}
\usepackage{booktabs}
\usepackage{longtable}
\usepackage{tabularx}
\usepackage{enumitem}
\usepackage{amssymb}
\usepackage{array}
\usepackage{parskip}
\usepackage{caption}
\usepackage{float}
\usepackage[destlabel]{hyperref}
\hypersetup{
    colorlinks,
    citecolor=green,
    filecolor=magenta,
    linkcolor=red,
    urlcolor=cyan
}
\setlength{\emergencystretch}{5em}


% ----- Numbered system tests (same list style as the SRS) -------------------
% Numbering is shared by every testlist.  Moving or inserting a test changes
% its displayed number and all \testref links automatically.
\makeatletter
\newlength{\reqlabelwidth}
\newlength{\reqlabelsep}
\setlength{\reqlabelsep}{0.75em}
\newsavebox{\reqlabelbox}
\newcommand{\reqsetwidth}[2]{\expandafter\gdef\csname reqwidth@#1\endcsname{#2}}
\newcommand{\reqlabel}[1]{%
  \sbox\reqlabelbox{\textbf{#1}}%
  \ifdim\wd\reqlabelbox>\csname reqmax@\reqkind\endcsname\relax
    \expandafter\xdef\csname reqmax@\reqkind\endcsname{\the\wd\reqlabelbox}%
  \fi
  \usebox\reqlabelbox}
\newenvironment{reqlist}[1]{%
  \def\reqkind{#1}%
  \@ifundefined{reqmax@#1}{\expandafter\gdef\csname reqmax@#1\endcsname{0pt}}{}%
  \@ifundefined{reqwidth@#1}%
    {\settowidth{\reqlabelwidth}{\textbf{ST-NFR-00:}}}%
    {\setlength{\reqlabelwidth}{\csname reqwidth@#1\endcsname}}%
  \begin{itemize}[leftmargin=\dimexpr\reqlabelwidth+\reqlabelsep\relax,
    labelwidth=\reqlabelwidth,labelsep=\reqlabelsep,
    align=left,itemsep=0.9em,parsep=0.35em]}%
  {\end{itemize}%
  \immediate\write\@auxout{\string\reqsetwidth{\reqkind}%
    {\csname reqmax@\reqkind\endcsname}}}
\makeatother

\newcounter{vnvtestnum}
\newenvironment{testlist}{\begin{reqlist}{STNFR}}{\end{reqlist}}
\newcommand{\testdef}[1]{%
  \refstepcounter{vnvtestnum}\label{test:#1}%
  \reqlabel{ST-NFR-\arabic{vnvtestnum}:}}
\newcommand{\testref}[1]{\hyperref[test:#1]{ST-NFR-\ref*{test:#1}}}

% ----- Usability Survey Numbering -----

\newcounter{surveyq}
\newcommand{\surveyq}{\stepcounter{surveyq}SQ\arabic{surveyq}}

% ----- Functional System Test Numbering -----

\newcounter{vnvfrtestnum}

\newenvironment{frtestlist}
    {\begin{reqlist}{STFR}}
    {\end{reqlist}}

\newcommand{\frtestdef}[1]{%
    \refstepcounter{vnvfrtestnum}%
    \label{test:#1}%
    \reqlabel{ST-FR-\arabic{vnvfrtestnum}:}}

\newcommand{\frtestref}[1]{%
    \hyperref[test:#1]{ST-FR-\ref*{test:#1}}}

% As in the Hazard Analysis, published SRS links use the label as their named
% destination; display numbers are written explicitly because the documents
% are compiled separately. This URL matches the SRS link used by the HA.
\newcommand{\srsurl}{https://team-hindsight.github.io/ophthamology-game/SRS/SRS.pdf}

% Functional requirement references
\newcommand{\srsfrref}[2]{\href{\srsurl\#nameddest=#1}{FR#2}}

% Non-functional requirement references
\newcommand{\srsnfrref}[2]{\href{\srsurl\#nameddest=#1}{NFR#2}}

% Links to other labelled SRS items, such as D3, T3, UC9, or LC4.
\newcommand{\srsitemref}[2]{\href{\srsurl\#nameddest=#1}{#2}}
\newcommand{\srsdoc}[1]{\href{\srsurl}{#1}}
\newcommand{\srstableref}[2]{\href{\srsurl\#nameddest=#1}{#2}}

% Local auxiliary-parameter entries are individually anchored: a reference
% jumps to the parameter's row, even in a long table spanning several pages.
\newcommand{\constdef}[1]{\hypertarget{const:#1}{\texttt{\detokenize{#1}}}}
\newcommand{\constref}[1]{\hyperlink{const:#1}{\texttt{\detokenize{#1}}}}

\input{../Comments.text}
\input{../Common.text}

\begin{document}

\title{System Verification and Validation Plan for \progname{}} 
\author{\authname}
\date{\today}
	
\maketitle

\pagenumbering{roman}

\begin{table}[!htbp]
\caption{Revision History} \label{TblRevisionHistory}
\begin{tabularx}{\textwidth}{llX}
\toprule
\textbf{Date} & \textbf{Developer(s)} & \textbf{Change}\\
\midrule
2026/10/30 & Mohammad Bilal & Completed section 3.3, individual reflection, \& contributed to sections 1, 4.2, 6.1, 6.2, 7 \\
2026/10/31 & Ahsan Muzammil & Completed section 2, 3.1, individual reflection, \& contributed to sections 1, 4.1, 4.2, 6.4, 6.5, 6.6, 7 \\
2026/10/31 & Jake Finlay & Completed section 3.6 \\
2026/10/31 & Jake Finlay & Completed the Rewards and Scoreboard tests in section 4.1 and their test data in the appendix \\
2026/11/01 & Jake Finlay & Completed the Look and Feel and the Usability and Humanity tests in section 4.2 \\
2026/11/01 & Jake Finlay & Completed individual reflection \\
2026/11/02 & Jake Finlay & Completed the usability survey questions in the appendix \\
\bottomrule
\end{tabularx}
\end{table}

\newpage

\tableofcontents

\newpage

\listoftables

\listoffigures
\wss{Remove this section if it isn't needed}

\newpage

\section{Symbols, Abbreviations, and Acronyms}

\renewcommand{\arraystretch}{1.2}
\begin{longtable}{@{}p{0.24\textwidth}p{0.69\textwidth}@{}}
\caption{Symbols, Abbreviations, and Acronyms}\label{tab:VnVAcronyms}\\
\toprule
\textbf{Symbol} & \textbf{Description}\\
\midrule
\endfirsthead
\multicolumn{2}{c}{\tablename\ \thetable{} -- continued from previous page}\\
\toprule
\textbf{Symbol} & \textbf{Description}\\
\midrule
\endhead
\midrule
\multicolumn{2}{r}{Continued on next page}\\
\endfoot
\bottomrule
\endlastfoot
AI & Artificial Intelligence.\\
AV & Audiovisual; synchronized narration and animation.\\
CD & Custom Document; an additional course deliverable.\\
CPU & Central Processing Unit.\\
FPS & Frames per second.\\
FR & Functional Requirement.\\
ID & Identifier; used to identify screens, assets, or content.\\
LC & Likely Change, as defined in the SRS.\\
MG & Module Guide.\\
MIS & Module Interface Specification.\\
NFR & Non-Functional Requirement.\\
OS & Operating System.\\
REB & Research Ethics Board.\\
SBOM & Software Bill of Materials.\\
SRS & Software Requirements Specification.\\
ST-FR & System test for a functional requirement; numbered in order of appearance.\\
ST-NFR & System test for a non-functional requirement; numbered in order of appearance.\\
TA & Teaching Assistant.\\
VnV & Verification and Validation.\\
\end{longtable}

For the remaining requirements terminology, see the definitions and
abbreviations table in the SRS.

\newpage

\pagenumbering{arabic}

This document presents the Verification and Validation (VnV) plan for
\progname{}, a game that helps children aged 3 to 7 prepare for eye
appointments by introducing them to common examination procedures. It outlines
how the team will verify that the software meets the requirements in the SRS
and validate that it is suitable for its intended users. The document covers
the planned reviews, system and unit tests, and validation activities. The
appendices provide supporting checklists, questionnaires, and test parameters.

\section{General Information}\label{sec:section2}

\subsection{Summary}

\progname{} is a mobile application designed to prepare children aged 3--7
for paediatric ophthalmology appointments. The application aims to reduce
children's anxiety and improve their preparedness by introducing common
eye examination procedures through age-appropriate interactive games.

The application follows a journey-based structure consisting of three
preparation levels: Vision Testing, Eye-Drop Preparation, and Slit-Lamp
Preparation. Each level focuses on a specific ophthalmology procedure
and contains smaller interactive activities or exercises that familiarize
children with what to expect during their appointments. Children can
access and replay any preparation level at any time, regardless of their
progress in other levels.

\progname{} provides visual demonstrations, narration, and encouraging
feedback throughout gameplay. Children can select and customize a
character, earn journey points and in-game rewards upon first completing
each preparation level, and track their progress through a personal
scoreboard. Completion status and rewards are tracked at the preparation
level rather than for individual exercises.

Caregivers can access procedure guides, manage application preferences,
and reset saved progress. The application operates offline after
installation, stores progress locally, and does not require user
accounts. \progname{} is intended solely as an educational preparation
tool and does not provide medical advice, diagnosis, or treatment.

\subsection{Objectives}\label{sec:Objectives}

The primary objective of this Verification and Validation (VnV) Plan
is to build confidence that \progname{} meets its specified requirements
and provides a safe, reliable, and age-appropriate preparation
experience for children aged 3--7. Given the project's scope and
available resources, verification and validation efforts will
prioritize the following qualities:

\begin{itemize}
    \item \textbf{Accuracy and Safety of Educational Content:}
    Verify that the preparation activities accurately represent
    common ophthalmology procedures, provide age-appropriate
    explanations, and avoid misleading or distressing information.
    Clinical content and applicable requirements will be reviewed
    with the project supervisor where feasible.

    \item \textbf{Functional Correctness and Reliability:}
    Verify that the application's core functionality behaves
    according to the SRS \cite{SRS}. Particular attention will be given to
    ensuring that all preparation levels remain accessible,
    gameplay interactions function correctly, progress and
    rewards are recorded accurately, and the application operates
    reliably offline.

    \item \textbf{Preliminary Usability and Accessibility:}
    Assess whether the application's navigation, instructions,
    and interactions are understandable and appropriate for
    the intended age group. This will involve reviews by the
    development team and peers, including task-based inspections
    in which reviewers evaluate the application from the
    perspectives of the child and caregiver personas defined
    in the SRS \cite{SRS}.
\end{itemize}

The following validation objectives are outside the scope of the
capstone project:

\begin{itemize}
    \item \textbf{Direct Usability Testing with Primary Users:}
    Usability testing involving children and their caregivers
    will not be conducted during the capstone project, as
    the required research ethics approval will not be pursued
    within the project timeline. Instead, preliminary usability
    assessments will be performed through peer reviews and
    persona based inspections. Since the reviewers are not the intended users, their feedback
    may not accurately reflect how children and caregivers would
    interact with the application. Therefore, these assessments
    can help identify potential usability issues but cannot fully
    confirm the application's usability for its intended users.

    \item \textbf{Clinical Effectiveness Validation:}
    Determining whether \progname{} reduces children's anxiety,
    improves cooperation during examinations, or increases
    clinical efficiency requires studies involving patients
    and healthcare professionals. Such studies are outside
    the scope of this project and may be pursued after the
    capstone, subject to research ethics approval.

    \item \textbf{Exhaustive Compatibility and Third-Party
    Verification:}
    Testing will focus on the devices and operating systems
    available to the team. Exhaustive testing across all
    supported device configurations and independent verification
    of third party software will not be performed due to
    resource constraints. Third party components will instead
    be assessed through their integration with the application.
\end{itemize}

The team will prioritize verification through software testing,
code and document reviews, and feedback from the supervisor where
appropriate. Reviews conducted by classmates using the child and caregiver
personas will help identify potential usability issues. However,
these reviews cannot fully confirm how actual children and
caregivers would interact with the application. The team will
document these limitations, while usability testing with the
intended users and evaluating the application's effectiveness
in reducing anxiety will be left for future work.

\subsection{Custom Documents}


\subsection{Relevant Documentation}

The following project documents provide the requirements,
design information, and development context used to plan
and conduct verification and validation activities for
\progname{}.

\begin{itemize}
    \item \textbf{Software Requirements Specification (SRS):}
    The SRS \cite{SRS} defines the functional and
    non-functional requirements of \progname{}, including
    preparation activities, journey progression, accessibility,
    and reliability. It serves as the primary reference for
    developing test cases and determining whether the
    application meets its specified requirements.

    \item \textbf{Problem Statement and Goals:}
    The Problem Statement \cite{ProblemStatement} identifies the problem
    \progname{} aims to address, its intended users, and the
    project's overall goals. It provides the context for
    evaluating whether the application supports its intended
    purpose, while recognizing that measuring actual anxiety
    reduction and clinical effectiveness is outside the
    scope of this capstone project.

    \item \textbf{Development Plan:}
    The Development Plan \cite{DevPlan} describes the team's
    development workflow, responsibilities, tools, and
    project schedule. It helps ensure that the planned
    verification and validation activities are feasible
    within the available time and resources.

    \item \textbf{Hazard Analysis:}
    The Hazard Analysis \cite{HazardAnalysis} identifies potential
    risks associated with the application and defines
    corresponding safety and security requirements.
    These findings will guide verification activities
    related to safe interactions, accurate educational
    content, data handling, and reliable application
    behaviour.

    \item \textbf{Module Guide (MG):}
    The Module Guide \cite{MG} describes the high level
    architecture of \progname{}, including its main modules,
    their responsibilities, and the relationships between
    them. It serves as a reference for verifying that the
    system architecture supports the functional and
    non functional requirements defined in the SRS.
    The MG will also be used during design verification
    to assess the organization of modules, their
    dependencies, and the overall structure of the system.

    \item \textbf{Module Interface Specification (MIS):}
    The Module Interface Specification \cite{MIS} provides
    the detailed design of \progname{}, including module
    interfaces, inputs, outputs, and expected behaviour.
    It will serve as the primary reference for planning
    unit tests and verifying that individual modules
    implement their specified functionality correctly.
    The MIS will also be reviewed alongside the MG
    to ensure consistency between the high level
    architecture and detailed module design.

\end{itemize}

The selected Fall and Winter Custom Documents will also
contribute to the VnV process by documenting preliminary
interface and usability evaluations. Their findings will
help identify design and usability issues that can be
addressed during development. References to these
documents will be added once their selection and content
have been finalized.

\newpage
\section{Plan}

This section outlines how the Hindsight team will carry out verification
and validation throughout the capstone project. It identifies the team members
responsible for these activities and explains how the SRS, design documents,
VnV plan, and implementation will be reviewed. It also describes the testing
tools and how the game will be validated with its intended users.

\subsection{Verification and Validation Team}\label{sec:vnvTeam}

Verification and validation activities for \progname{} will
be shared among the five members of Team 2, Hindsight.
Each member has been assigned functional and nonfunctional
requirement categories for testing based on the SRS
\cite{SRS}. These assignments are informed by the
team's development roles \cite{DevPlan}, but verification
will remain a collaborative effort. 

Each member will coordinate the development and execution
of test cases for their assigned requirement categories,
document testing results, and report identified defects.
However, these assignments do not imply exclusive
responsibility for verifying a particular part of the
application. Team members will also review one another's
test cases, implementation, and results where feasible. 

Responsibilities may vary between SRS verification, design
verification, and implementation verification to encourage
independent assessment. The project supervisor will provide
feedback on selected requirements and educational content
where appropriate. Classmates may also participate in
preliminary usability evaluations using the child and
caregiver personas defined in the SRS. 

Table~\ref{tab:VnVTeam} summarizes each member's
development role, assigned testing categories, and
associated verification responsibilities. 



\begin{table}[H]
\centering
\footnotesize
\setlength{\tabcolsep}{3pt}
\renewcommand{\arraystretch}{1.15}

\caption{VnV Team Roles and Testing Responsibilities}
\label{tab:VnVTeam}

\begin{tabularx}{\textwidth}{
    @{}
    >{\raggedright\arraybackslash}p{2.4cm}
    >{\raggedright\arraybackslash}p{2.3cm}
    >{\raggedright\arraybackslash}p{5.2cm}
    >{\raggedright\arraybackslash}X
    @{}
}
\toprule
\textbf{Role} &
\textbf{Member Name} &
\textbf{Testing Areas (FR/NFR)} &
\textbf{VnV Responsibilities} \\
\midrule

Technical Manager &
Ahsan Muzammil &
\textbf{FR:}
Setup~(\ref{sec:AMSetupFRTests}),
Journey Map~(\ref{sec:AMJourneyFRTests}),
Level Content~(\ref{sec:AMLevelContentFRTests}).
\newline
\textbf{NFR:}
Cultural and Political~(\ref{sec:AMCulturalNFRTests}). &
Test setup, level accessibility, and procedure content.
Coordinate testing and defect tracking. \\

\addlinespace[4pt]

Lead Front end
Developer &
Jake Finlay &
\textbf{FR:}
Rewards and Scoreboard~(\ref{sec:RewardsScoreboardFRTests}).
\newline
\textbf{NFR:}
Look and Feel~(\ref{sec:LookAndFeelNFRTests}),
Usability and Humanity~(\ref{sec:UsabilityHumanityNFRTests}). &
Test interface presentation, rewards,
and usability requirements. \\

\addlinespace[4pt]

Lead Back end
Developer &
Mohammad Bilal &
\textbf{NFR:}
Performance~(\ref{sec:PerformanceNFRTests}),
Operational and Environmental~(\ref{sec:OperationalNFRTests}),
Maintainability and Support~(\ref{sec:MaintainabilityNFRTests}),
Security~(\ref{sec:SecurityNFRTests}). &
Test performance, reliability,
maintainability, and security. \\

\addlinespace[4pt]

Lead Researcher &
Jackson Beach &
\textbf{FR:}
Instructions and Feedback~(\ref{sec:InstructionsFeedbackFRTests}),
Exiting and Interruptions~(\ref{sec:ExitingInterruptionsFRTests}),
Saving and Restoring Progress~(\ref{sec:SavingRestoringFRTests}). &
Test instructional behaviour,
interruption handling, and
progress recovery. \\

\addlinespace[4pt]

Team Liaison &
Ethan McMehen &
\textbf{FR:}
Characters and Accessories~(\ref{sec:CharactersAccessoriesFRTests}),
Caregiver Features~(\ref{sec:CaregiverFeaturesFRTests}).
\newline
\textbf{NFR:}
Compliance~(\ref{sec:ComplianceNFRTests}). &
Test customization, caregiver
features, and compliance.
Coordinate external feedback. \\

\addlinespace[4pt]

Clinical Reviewer &
Project Supervisor &
\textbf{Clinical Review:}
Selected requirements and
educational content. &
Provide feedback on clinical accuracy,
age-appropriateness, and procedure
representations where feasible. \\

\bottomrule
\end{tabularx}
\end{table}


\subsection{SRS Verification}

\wss{List any approaches you intend to use for SRS verification.  This may
  include ad hoc feedback from reviewers, like your classmates (like your
  primary reviewer), or you may plan for something more rigorous/systematic.}

\wss{If you have a supervisor for the project, you shouldn't just say they will
read over the SRS.  You should explain your structured approach to the review.
Will you have a meeting?  What will you present?  What questions will you ask?
Will you give them instructions for a task-based inspection?  Will you use your
issue tracker?}

\wss{Maybe create an SRS checklist?}

\subsection{Design Verification}\label{sec:DesignVerification}

Design verification will be carried out after the Module Interface Specification
(MIS) is completed. The purpose is to check that the architecture described in
the Module Guide (MG) and the detailed design in the MIS are consistent with
the requirements in the SRS and address the relevant concerns identified
in the Hazard Analysis. The design will be verified through three reviews:

\begin{enumerate}
    \item \textbf{Structured Internal Review:} All members of the Hindsight team
    will review the MG and MIS, preferably in a group meeting. One team member
    will be appointed to lead the review and work through the checklist in
    Appendix~\ref{sec:DesignVerificationChecklist}. The team will discuss any
    concerns, record them as GitHub issues, and assign members to address them.

    \item \textbf{External Peer Review:} Members of at least one other capstone
    team will review the design as part of the course's milestone peer review
    process. The checklist in Appendix~\ref{sec:DesignVerificationChecklist}
    will be shared with the reviewers, but they will not be required to use it.
    Their findings will be submitted as GitHub issues for the team to review
    and address.

    \item \textbf{Teaching Assistant Feedback:} The teaching assistant (TA)
    will evaluate the design documents as part of the course deliverable
    assessment. Any feedback requiring changes will be recorded as GitHub
    issues and addressed by the team.
\end{enumerate}

For each review, GitHub issues will identify the relevant section of the design,
the problem found, and the change needed. The assigned team member will make
the correction, and another team member will check the update before the issue
is closed.

\subsection{Verification and Validation Plan Verification}

The Verification and Validation Plan will be reviewed by the team before submission and after major
revisions to check completeness, accuracy, and feasibility. Each section
will be checked by a team member besides its author. Colleagues from at
least one other team will also review the plan through the peer review process.
This checklist will guide the review:

\begin{itemize}[label={$\square$}, leftmargin=*, itemsep=0.35em]
    \item Are all sections complete, with any deferred or
    excluded work explained?

    \item Does every functional and non-functional requirement have a planned
    test or review, and does the traceability table match the tests?

    \item Does each test specify its initial conditions, inputs, steps, expected results, and evidence?

    \item Are expected results and test parameters consistent with the SRS
    and Hazard Analysis (and with the MG and MIS when available)?

    \item Are testing, reviews, and other validation activities feasible with the available time, equipment, and participants?

    \item Are terminology, test identifiers, and references consistent and
    correct?
\end{itemize}

Feedback will be tracked through GitHub issues. Each issue will identify the affected section, the problem,
and the proposed change, and will be assigned to a team member. Corrections
will be committed with reference to the issue and checked by another member
before it is closed. If a suggestion is not adopted, the team will record the
reason why in the issue.

\subsection{Implementation Verification}

\wss{You should at least point to the tests listed in this document and the unit
  testing plan.}

\wss{In this section you would also give any details of any plans for static
  verification of the implementation.  Potential techniques include code
  walkthroughs, code inspection, static analyzers, etc.}

\wss{The final class presentation in CAS 741 could be used as a code
walkthrough.  There is also a possibility of using the final presentation (in
CAS741) for a partial usability survey.}

\subsection{Automated Testing and Verification Tools}\label{sec:AutomatedTools}

\progname{} is written in C\# with the Unity engine. The Development Plan
\cite{DevPlan} already sets out the main tool choices. Section~8.3 describes the
GitHub Actions pipeline, Section~11 lists the unit testing framework, and
Section~13 gives the coding standard and how it is enforced. This section
summarizes the tools that automate testing, formatting, and linting, and adds
the plan for code coverage. Table~\ref{tab:AutomatedTools} lists each tool.
These choices may change as implementation progresses. Any change will be
recorded in the revision history of this document and the Development Plan.

\begin{table}[!htbp]
\centering
\caption{Automated Testing and Verification Tools}\label{tab:AutomatedTools}
\renewcommand{\arraystretch}{1.2}
\begin{tabularx}{\textwidth}{@{}p{0.19\textwidth}p{0.29\textwidth}X@{}}
\toprule
\textbf{Purpose} & \textbf{Tool} & \textbf{Use}\\
\midrule
Unit testing & Unity Test Framework (NUnit), Edit Mode tests &
  Tests game logic that does not need a running scene, such as level
  progression, quiz scoring, and saved progress.\\
Unity interaction testing & Unity Test Framework, Play Mode tests &
  Tests scene loading, touch input handling, and component behaviour across
  frames.\\
Continuous integration & GitHub Actions with GameCI (\texttt{unity-test-runner},
  \texttt{unity-builder}) &
  Runs the format check, linting, both test suites, and a build on every pull
  request. All checks must pass before a merge (Development Plan,
  Section~8.3).\\
Formatting & \texttt{dotnet format} with a shared \texttt{.editorconfig} &
  Fails CI when code breaks the formatting or naming rules in Development Plan
  Section~13.\\
Linting & .NET Roslyn analyzers and Microsoft.Unity.Analyzers &
  Reports common C\# and Unity mistakes, such as unused code and inefficient
  Unity API calls. Rider inspections give the same feedback locally.\\
Code coverage & Unity Code Coverage package &
  Records line and method coverage during CI test runs.\\
\bottomrule
\end{tabularx}
\end{table}

\textbf{Formatting and Linting}

The naming rules in the Development Plan,
such as PascalCase methods and camelCase private fields, will be written as
\texttt{.editorconfig} rules. CI will run \texttt{dotnet format} in
verification mode against the project files that Unity generates, so any
violation fails the pull request. The analyzer rules will be configured in the
same \texttt{.editorconfig}. Rules that catch real defects will be set to
errors and fail the build. Style suggestions will stay as warnings. Reviewers
will still check the parts of the standard that tools cannot, such as comment
quality and separation of game logic from UI logic.

\textbf{Code Coverage}

GameCI will run the Edit Mode and Play Mode tests
with the Unity Code Coverage package enabled. Coverage will be measured only
for the team's game assemblies. Test code, Unity packages, and Asset Store
content will be excluded. Each CI run will produce an HTML report and a
summary of line and method coverage for each assembly. The summary will be
shown on the pull request and the full report saved as a workflow artifact.
For each milestone, the team will record the coverage of each assembly in the
VnV Report and list any modules with little or no coverage. The target is
\constref{COVERAGE_LINE_TARGET} line coverage for game-logic assemblies. UI and
scene code will rely more on Play Mode and manual system tests. Coverage will
be used to find untested code, not as proof of correctness.

Performance is not measured in the pipeline. The Unity Profiler is used
manually during the performance system tests, such as \testref{ST_Smoothness}
and \testref{ST_DeviceHeat}.

\subsection{Software Validation}

\wss{If there is any external data that can be used for validation, you should
  point to it here.  If there are no plans for validation, you should state that
  here.}

\wss{You might want to use review sessions with the stakeholder to check that
the requirements document captures the right requirements.  Maybe task based
inspection?}

\wss{For those capstone teams with an external supervisor, the Rev 0 demo should 
be used as an opportunity to validate the requirements.  You should plan on 
demonstrating your project to your supervisor shortly after the scheduled Rev 0 demo.  
The feedback from your supervisor will be very useful for improving your project.}

\wss{For teams without an external supervisor, user testing can serve the same purpose 
as a Rev 0 demo for the supervisor.}

\wss{This section might reference back to the SRS verification section.}

\newpage
\section{System Tests}

This section describes the system tests used to verify the functional and
non-functional requirements in the SRS. The functional tests check that
the main features of the game work as expected, while the non-functional tests
check qualities such as performance, offline use, maintainability, and security.
Each test explains what will be checked and how the result will be recorded.

\subsection{Tests for Functional Requirements}

The tests in this section verify the functional requirements (FRs)
defined in the SRS. These requirements describe the features and
behaviours the game must provide. The tests are organized according
to the FR categories in the SRS, covering setup, journey map,
level content, instructions and feedback, exiting and interruptions,
saving and restoring progress, rewards and scoreboard, characters
and accessories, and caregiver features. Each test identifies the
requirements it covers, the initial conditions, expected results,
and how the test will be performed.

For the Capstone course, functional testing will primarily be
performed manually using the available test devices. Testers will
record the device, operating system, build version, and test
results. Any limitations or changes to the planned testing
procedures will be documented under the relevant test.

\subsubsection{Setup}\label{sec:AMSetupFRTests}

These tests cover the first-use medical notice and the locally saved
caregiver acknowledgement in \srsfrref{FR_MedicalNotice}{1}--\srsfrref{FR_RecordAck}{2} of the SRS \cite{SRS}.
For tests requiring a first launch, use a test installation with its
application data cleared and no restored backup. Record the device,
operating-system version, build identifier, and observed result.

\begin{frtestlist}

\item[\frtestdef{ST_AM_NoticeFirstLaunch}]\textbf{Medical Notice Before First Use}

\textit{Type:} Functional, Dynamic, manual.

\textit{Requirements Covered:} \srsfrref{FR_MedicalNotice}{1}.

\textit{Initial State:} The current build is installed on a test device.
All application data has been cleared, and no caregiver acknowledgement
is stored.

\textit{Input/Condition:} Launch the app. Read the medical notice and
attempt to enter the home screen without selecting its acknowledgement
control. Close and reopen the app without acknowledging the notice.

\textit{Output/Result:} The notice identifies the game as an educational
preparation aid and states that it does not replace professional medical
advice, diagnosis, or treatment. The home screen remains inaccessible
until acknowledgement. The notice is shown again after relaunch.

\textit{Test Case Derivation:} \srsfrref{FR_MedicalNotice}{1} requires the notice before first use
and prevents access to the home screen until the caregiver acknowledges it.

\textit{How test will be performed:} On a clean installation, capture the
notice and check all required statements. Attempt any available back,
home, or close-navigation action without acknowledging. Relaunch and
confirm the notice is still required. Record the available controls and
whether any route bypasses acknowledgement.

\item[\frtestdef{ST_AM_NoticeInGuide}]\textbf{Medical Notice in Caregiver Guide}

\textit{Type:} Functional, Dynamic, manual.

\textit{Requirements Covered:} \srsfrref{FR_MedicalNotice}{1}.

\textit{Initial State:} Initial setup has been completed and the home
screen is accessible.

\textit{Input/Condition:} Open the caregiver procedure guide and select
the available medical notice or educational-use disclaimer.

\textit{Output/Result:} The medical notice remains accessible and states
that the game is for educational preparation, not medical advice,
diagnosis, or treatment. Viewing it does not change journey progress.

\textit{Test Case Derivation:} \srsfrref{FR_MedicalNotice}{1} explicitly requires the same notice
to remain available from the caregiver procedure guide.

\textit{How test will be performed:} Navigate to the caregiver guide,
open the notice, compare its required statements with those checked in
\textit{ST\_AM\_NoticeFirstLaunch}, and return home. Record the
navigation path and capture the displayed notice.

\item[\frtestdef{ST_AM_AcknowledgementSaved}]\textbf{Acknowledgement Persists Across Launches}

\textit{Type:} Functional, Dynamic, manual.

\textit{Requirements Covered:} \srsfrref{FR_RecordAck}{2}.

\textit{Initial State:} A clean installation is displaying the first-use
medical notice; no acknowledgement has been recorded.

\textit{Input/Condition:} Select the acknowledgement control once,
complete any introductory screens, reach the home screen, fully close
the app, and launch it again. Repeat the close-and-launch operation once.

\textit{Output/Result:} The acknowledgement is saved locally. Both
subsequent launches reach the home screen without requesting acknowledgement
again. No account or network connection is required to retain it.

\textit{Test Case Derivation:} \srsfrref{FR_RecordAck}{2} requires a locally stored acknowledgement
and omission of the notice on later launches.

\textit{How test will be performed:} Record the acknowledgement action,
then force-close and relaunch the app twice. Confirm that the home screen
is reached without a repeated notice. If a development build exposes
the acknowledgement field, also check that it changes from false to true.

\end{frtestlist}

\subsubsection{Journey Map}\label{sec:AMJourneyFRTests}

These tests verify \srsfrref{FR_JourneyMap}{3}--\srsfrref{FR_Replay}{5}: presentation of the three procedure levels,
independent access to every level, and replay without loss of a completed
status. The names, order, and allowed statuses come from the SRS
\cite{SRS}; no custom image or data file is needed because the test
uses the levels and illustrations packaged in the release build.

\begin{frtestlist}

\item[\frtestdef{ST_AM_JourneyFresh}]\textbf{Journey Map on a New Profile}

\textit{Type:} Functional, Dynamic, manual.

\textit{Requirements Covered:} \srsfrref{FR_JourneyMap}{3}.

\textit{Initial State:} Initial setup is acknowledged, and a new
journey has zero completed or started levels. If necessary, use the
caregiver's journey-reset control before starting the test.

\textit{Input/Condition:} From the home screen, open the journey map.
Do not start any level.

\textit{Output/Result:} Exactly the three defined procedure levels are
shown in level-number order: Vision Testing, Eye-Drop Preparation, and
Slit-Lamp Preparation. Each has a name, an illustration, and a visible
\textit{Not Started} status. No level is displayed as locked.

\textit{Test Case Derivation:} \srsfrref{FR_JourneyMap}{3} specifies one illustrated, named,
status-labelled map entry per procedure, ordered by level number.

\textit{How test will be performed:} Reset journey progress, open the
map, and record the order, titles, images, status labels, and availability
of all three entries in a screenshot and checklist.

\item[\frtestdef{ST_AM_JourneyMixedStatuses}]\textbf{Map Displays Three Different Statuses}

\textit{Type:} Functional, Dynamic, manual.

\textit{Requirements Covered:} \srsfrref{FR_JourneyMap}{3},
\srsfrref{FR_StartLevel}{4}, \srsfrref{FR_LevelEnd}{11}.

\textit{Initial State:} Initial setup is complete; journey progress is
reset so all three levels are \textit{Not Started}.

\textit{Input/Condition:} Start Vision Testing, then use its home
control before the ending. Next start Eye-Drop Preparation and reach
its end-of-level screen. Do not open Slit-Lamp Preparation. Return to
the journey map.

\textit{Output/Result:} Vision Testing is \textit{In Progress},
Eye-Drop Preparation is \textit{Completed}, and Slit-Lamp Preparation
is \textit{Not Started}. All three entries remain selectable.

\textit{Test Case Derivation:} \srsfrref{FR_JourneyMap}{3} requires the map to show the current
status of each level; FR4 and FR11 define the changes on start and finish.

\textit{How test will be performed:} Follow the stated sequence on a
reset journey, recording each transition. Compare all three status
indicators with the expected state and capture the final journey map.

\item[\frtestdef{ST_AM_IndependentLevelAccess}]\textbf{Start Any Level Without Unlocking}

\textit{Type:} Functional, Dynamic, manual.

\textit{Requirements Covered:} \srsfrref{FR_StartLevel}{4}.

\textit{Initial State:} Journey progress is reset; all levels are
\textit{Not Started}, and journey points and coins are zero.

\textit{Input/Condition:} Select Slit-Lamp Preparation first and
activate its play control. Before the ending, return to the map.
Then select Eye-Drop Preparation, activate play, and return early.
Finally select Vision Testing and activate play. For each start, record
the first displayed level step.

\textit{Output/Result:} Each selected level starts at its beginning
without requiring a completed prerequisite, journey points, or coins.
Its status becomes \textit{In Progress}; the other levels retain their
existing statuses. No reward is given merely for starting a level.

\textit{Test Case Derivation:} \srsfrref{FR_StartLevel}{4} allows every level to start regardless
of progress or rewards, applies the start operation, and requires the
level to begin from its first step.

\textit{How test will be performed:} Execute the three starts in the
specified order. After each early exit, inspect the journey map and
record the starting scene, all three statuses, and displayed reward
balances. Flag any lock, prerequisite, or resumed mid-level scene.

\item[\frtestdef{ST_AM_ReplayCompleted}]\textbf{Replay Does Not Remove Completion}

\textit{Type:} Functional, Dynamic, manual.

\textit{Requirements Covered:} \srsfrref{FR_Replay}{5},
\srsfrref{FR_StartLevel}{4}.

\textit{Initial State:} On a reset journey, Vision Testing has been
played to its end once and shows \textit{Completed}. Record the
current journey points, coins, and level status.

\textit{Input/Condition:} Select the completed Vision Testing level,
start its replay, exit before the ending, and inspect the map. Start
Vision Testing again and play to the end.

\textit{Output/Result:} Both replays start from the beginning.
Vision Testing remains \textit{Completed} during and after replay,
including after early exit. Its original completion reward is not
awarded again; recorded points and coins remain unchanged by replay.

\textit{Test Case Derivation:} \srsfrref{FR_Replay}{5} preserves completed status during
replay, while FR4 restarts the level from the beginning. The SRS
formal progress model permits first-completion rewards only once.

\textit{How test will be performed:} Capture the map and reward balances
before replay, after early exit, and after replay completion. Compare
these values and record whether each replay began at the introduction.

\end{frtestlist}

\subsubsection{Level Content}\label{sec:AMLevelContentFRTests}

These tests verify \srsfrref{FR_LevelContent}{6}--\srsfrref{FR_LevelEnd}{11}
using the three levels and their narration, images, and animations.
Before testing, the tester will record the build version and prepare
a content inventory of relevant assets, scripts, age-group variants,
and procedure steps, while clinical accuracy will be reviewed
separately in the NFR tests.

\begin{frtestlist}

\item[\frtestdef{ST_AM_ProcedureContent}]\textbf{Procedure Content in Every Level}

\textit{Type:} Functional, Dynamic, manual content inspection.

\textit{Requirements Covered:} \srsfrref{FR_LevelContent}{6}.

\textit{Initial State:} Initial setup is complete and the three
procedure levels are available with the final content inventory.

\textit{Input/Condition:} Open Vision Testing, Eye-Drop Preparation,
and Slit-Lamp Preparation separately. For each, advance through all
built-in screens and animations, including any optional exercises.

\textit{Output/Result:} Each level presents at least one identifiable
step, piece of equipment, or instruction for its own procedure:
Vision Testing includes an eye-chart or symbol-identification activity;
Eye-Drop Preparation depicts or explains clinician-administered drops;
and Slit-Lamp Preparation depicts or explains the slit-lamp device or
patient positioning. None is an unrelated activity without procedure
information.

\textit{Test Case Derivation:} \srsfrref{FR_LevelContent}{6} requires each level to prepare for
its corresponding procedure by presenting at least one relevant part.
The three use cases describe the intended procedure examples.

\textit{How test will be performed:} Play each level from its beginning,
record the first qualifying procedure moment and its asset or screen ID,
and capture a screenshot or narrated-screen recording. Compare the
three observations with the approved procedure inventory.

\item[\frtestdef{ST_AM_AgeGroupSelection}]\textbf{Age-Specific Content Selection}

\textit{Type:} Functional, Dynamic, manual.

\textit{Requirements Covered:} \srsfrref{FR_AgeGroup}{7}.

\textit{Initial State:} The caregiver can access preferences. Before
testing, the content inventory states which levels, if any, have a
3--4-year variant, a 5--7-year variant, and default content.

\textit{Input/Condition:} Set the age group to 3--4 years and start
each level from the beginning; record its opening content ID. Repeat
with age group 5--7 years. Finally use the unset age-group option, or clear local preferences
and complete setup again if no unset control is exposed, then start
each level again.

\textit{Output/Result:} For each run, a level presents its matching
age-specific content when that variant exists. Where no matching
variant exists, or the age group is unset, it presents the listed
default content. No level becomes unavailable because of the setting.

\textit{Test Case Derivation:} \srsfrref{FR_AgeGroup}{7} defines the values \textit{3--4},
\textit{5--7}, and \textit{unset}, and requires a matching variant
when provided, with fallback to default otherwise.

\textit{How test will be performed:} Use the caregiver preferences to
run the stated three-setting matrix across all three levels. Compare
each displayed content ID with the recorded content inventory; record
all nine setting-and-level combinations. Do not mark a missing optional
age-specific variant as a failure if its default content is used.

\item[\frtestdef{ST_AM_SensationBeforeSkip}]\textbf{Required Sensation Explanation Cannot Be Skipped}

\textit{Type:} Functional, Dynamic, manual.

\textit{Requirements Covered:} \srsfrref{FR_DescribeSensations}{8}.

\textit{Initial State:} The current build includes the approved
procedure narration and animations. Identify every scene that depicts
a potentially felt sensation (for example, stinging from eye drops
or bright light during a slit-lamp examination) in the content inventory.

\textit{Input/Condition:} For each identified scene, start its level
from the beginning and use the available skip control as soon as it
appears. Repeat once without skipping the optional animation or exercise.

\textit{Output/Result:} On both paths, the child receives the relevant
sensation description before the sensation is depicted or the optional
part is skipped. The mandatory description itself cannot be bypassed.
No narration, text, or image promises that the procedure will be
painless or free of discomfort.

\textit{Test Case Derivation:} \srsfrref{FR_DescribeSensations}{8} requires honest advance
explanations of sensations and prohibits skipping the description,
even when the corresponding exercise or animation is skipped.

\textit{How test will be performed:} Record both paths for every scene
listed in the inventory. Note the first frame or narration cue for the
description, the skip action, and the first frame depicting the
sensation. Review all text and narration for prohibited reassurance.
Save the recordings and a pass/fail result for every applicable scene.

\item[\frtestdef{ST_AM_ProcedureVariation}]\textbf{Differences Between Game and Clinic}

\textit{Type:} Functional, Dynamic, manual.

\textit{Requirements Covered:} \srsfrref{FR_ClinicsDiffer}{9}.

\textit{Initial State:} All three levels contain the current approved
narration and pictures. Device audio is enabled.

\textit{Input/Condition:} Play each level from beginning to end,
including its introductory and concluding content.

\textit{Output/Result:} Each of the three levels communicates, through
both narration and pictures, that a real doctor may perform the
procedure somewhat differently from the game. The message is present
in Vision Testing, Eye-Drop Preparation, and Slit-Lamp Preparation.

\textit{Test Case Derivation:} \srsfrref{FR_ClinicsDiffer}{9} explicitly requires this message in
both narration and pictures in every level.

\textit{How test will be performed:} For each level, identify the
narration clip or on-screen spoken message and the corresponding
picture or visual cue. Record both content IDs and a screenshot or
recording. Mark a level as failing if either channel is missing.

\item[\frtestdef{ST_AM_ExpressConcern}]\textbf{Permission to Tell the Doctor About Discomfort}

\textit{Type:} Functional, Dynamic, manual.

\textit{Requirements Covered:} \srsfrref{FR_SpeakUp}{10}.

\textit{Initial State:} All three levels are accessible and audio
is enabled.

\textit{Input/Condition:} Complete each level once normally. Reset
the journey, then complete each level using skip for every optional
exercise or animation that offers it.

\textit{Output/Result:} Before the ending in both paths for all three
levels, the child is told that they may tell the doctor if they feel
worried or uncomfortable. Skipping exercises does not remove this message.

\textit{Test Case Derivation:} \srsfrref{FR_SpeakUp}{10} requires the reassurance before
the end of every level. The alternative skip path checks that a child
cannot miss it by bypassing an optional exercise.

\textit{How test will be performed:} Record the normal and skip-heavy
runs for all three levels, identify the screen or narration ID conveying
the message, and confirm that it occurs before each end-of-level screen.

\item[\frtestdef{ST_AM_LevelCompletion}]\textbf{Completion Recorded Only at Level Ending}

\textit{Type:} Functional, Dynamic, manual.

\textit{Requirements Covered:} \srsfrref{FR_LevelEnd}{11}.

\textit{Initial State:} Journey progress is reset so all three
levels are \textit{Not Started}.

\textit{Input/Condition:} Start Vision Testing and exit before its
ending. Return to the map. Then play Vision Testing to its end,
followed by Eye-Drop Preparation and Slit-Lamp Preparation to their
respective endings. Record the map after each step.

\textit{Output/Result:} Early exit from Vision Testing leaves it
\textit{In Progress} without an end-of-level screen or completion
reward. Reaching the ending of each level displays its end-of-level
screen and changes only that level's status to \textit{Completed}.
All three are \textit{Completed} only after their own endings have been
reached.

\textit{Test Case Derivation:} \srsfrref{FR_LevelEnd}{11} applies the finish operation at
level ending and changes the level status to \textit{Completed}. The
formal model does not count individual exercises as separate levels.

\textit{How test will be performed:} Follow the specified exit and
three-completion sequence. Capture the ending of each level and the
map after each exit or completion; compare all three statuses after
every transition.

\item[\frtestdef{ST_AM_SkippedActivitiesStillComplete}]\textbf{Skipped Exercises Do Not Prevent Level Completion}

\textit{Type:} Functional, Dynamic, manual.

\textit{Requirements Covered:} \srsfrref{FR_LevelEnd}{11},
\srsfrref{FR_DescribeSensations}{8},
\srsfrref{FR_SpeakUp}{10}.

\textit{Initial State:} Journey progress is reset, and the three
levels have their normal content and available skip controls.

\textit{Input/Condition:} Start each level. For every optional
exercise, round, or animation that offers skip, select skip instead of
completing it; continue through any mandatory educational narration and
reach the level ending.

\textit{Output/Result:} Each level reaches its end of level screen
and becomes \textit{Completed}, despite skipped optional activities.
Mandatory sensation descriptions and the message about telling a doctor
about worry or discomfort are still presented. No individual skipped
exercise is incorrectly treated as a separate unfinished level.

\textit{Test Case Derivation:} The SRS tracks completion per
procedure level, not per exercise; \srsfrref{FR_LevelEnd}{11} requires completion at the
ending. \srsfrref{FR_DescribeSensations}{8} and \srsfrref{FR_SpeakUp}{10} protect mandatory educational information.

\textit{How test will be performed:} Record a skip heavy run through
all three levels. Count and note each selected skip, verify every
mandatory message, and capture each end of level screen and final
journey map status.

\end{frtestlist}


% ============================================================
% 4. Instructions and Feedback
% ============================================================

\subsubsection{Instructions and Feedback}
\label{sec:InstructionsFeedbackFRTests}

These tests verify \srsfrref{FR_SpokenInstructions}{12}-\srsfrref{FR_IgnoreRepeatedTouches}{19}.
Extend the content inventory from Section~\ref{sec:AMLevelContentFRTests}
with child-facing instructions, exercises, questions and reviewed answers,
animations, and screen states and transitions. Include instructions outside
levels and check all levels' default content and each distinct
age-group variant.

Use walkthroughs and the Unity test tools in the Development Plan, with
temporary input, playback, and storage logs where needed. Record the screen
or content ID with supporting recordings or logs.

\begin{frtestlist}

\item[\frtestdef{ST_InstructionsFeedback_01}]
\textbf{Narrated Instructions with Demonstrations}

\textit{Type:} Functional, Dynamic, manual.

\textit{Requirements Covered:} \srsfrref{FR_SpokenInstructions}{12}.

\textit{Initial State:} Setup is complete and narration is audible.

\textit{Input/Condition:} Follow every listed child-facing instruction
using only narration and demonstrations. Include hints, answer feedback,
and instructions outside levels.

\textit{Output/Result:} Every instruction has matching narration and a
visual demonstration. No required action depends on reading.

\textit{Test Case Derivation:} \srsfrref{FR_SpokenInstructions}{12} requires both forms of guidance and
prohibits levels from requiring reading.

\textit{How test will be performed:} Record the walkthrough with audio.
Match each instruction ID to its narration, demonstration, and action;
note missing or conflicting guidance and any reliance on written text.

\item[\frtestdef{ST_InstructionsFeedback_02}]
\textbf{Hints after Incorrect Actions}

\textit{Type:} Functional, Dynamic, manual.

\textit{Requirements Covered:} \srsfrref{FR_HintOnMiss}{13}.

\textit{Initial State:} An exercise awaits an action identified in the
inventory. Wrong answers to questions are covered by
\frtestref{ST_InstructionsFeedback_05}.

\textit{Input/Condition:} Make an incorrect action, repeat it, then perform
the requested action. Cover each distinct incorrect response path, such as
placing or releasing an object outside its target. Unrelated touches are
checked in \frtestref{ST_ExitingInterruptions_06}.

\textit{Output/Result:} Each incorrect action gives an encouraging hint
and allows another attempt. The correct action is accepted and play continues.

\textit{Test Case Derivation:} \srsfrref{FR_HintOnMiss}{13} requires hints and retries for incorrect
actions, including repeated mistakes.

\textit{How test will be performed:} Record each action, the hint's wording
and visual guidance, and the successful retry. Check that the hint guides
without blame or discouragement.

\item[\frtestdef{ST_InstructionsFeedback_03}]
\textbf{Visual Cues Remain Available}

\textit{Type:} Functional, Dynamic, manual.

\textit{Requirements Covered:} \srsfrref{FR_VisualCues}{14}.

\textit{Initial State:} Setup is complete. Visual-cue size and reduced-motion
settings are available under \srsfrref{FR_Preferences}{41}.

\textit{Input/Condition:} Inspect every listed level state and transition,
including after skips and restarts. Try changing or disabling cues through
child-facing controls, and accessing cue settings before and after the
parental gate. Check minimum and maximum cue sizes and reduced motion on
and off, varying one setting at a time, then combining minimum size with
reduced motion.

\textit{Output/Result:} Cues remain visible and provide guidance in every
state. The child cannot disable them; settings that change them require
the parental gate. Size and motion changes do not remove needed guidance.

\textit{Test Case Derivation:} \srsfrref{FR_VisualCues}{14} requires persistent cues and restricts
changes to caregiver settings under \srsfrref{FR_ParentGate}{38}.

\textit{How test will be performed:} Mark each state and setting combination
in the inventory. Record missing cues and whether each route to cue settings
requires the gate.

\item[\frtestdef{ST_InstructionsFeedback_04}]
\textbf{Instructions without Sound}

\textit{Type:} Functional, Dynamic, manual.

\textit{Requirements Covered:} \srsfrref{FR_NarrationFallback}{15}.

\textit{Initial State:} Setup is complete. Confirm audible narration before
each sound condition and identify device mute modes that silence media,
not just calls or alerts.

\textit{Input/Condition:} Follow every listed instruction in a silent
walkthrough, including hints, reveals, and sensation descriptions. Exercise
each mute-detection and prompt path using device mute, zero media volume,
and zero caregiver narration volume. Cover silence at entry and loss of
sound during playback. Leave the prompt unanswered and finish the level
visually, then restore sound using the indicated control and play the next
instruction.

\textit{Output/Result:} Pictures, icons, and demonstrations convey every
instruction. The caregiver prompt explains how to restore sound without
blocking play. Muting does not pause play or prevent completion; restoring
sound permits subsequent narration while visual guidance remains available.

\textit{Test Case Derivation:} \srsfrref{FR_NarrationFallback}{15} requires visual guidance and a
non-blocking prompt. SR-18 of the Hazard Analysis \cite{HazardAnalysis}
also covers narration disabled through settings.

\textit{How test will be performed:} Verify silence at the device output.
Record the guidance, prompt, completion, and sound restoration. Record
unavailable mute modes as unverified.

\item[\frtestdef{ST_InstructionsFeedback_05}]
\textbf{Correct Answers after a Wrong Answer}

\textit{Type:} Functional, Dynamic, manual with timing logs.

\textit{Requirements Covered:} \srsfrref{FR_RevealAnswer}{16},
\srsfrref{FR_NoPenalty}{18}; reveal cues and timing from
\srsnfrref{NFR_Colour}{10} and \srsnfrref{NFR_Responsiveness}{12}.

\textit{Initial State:} A question is displayed with audible narration
and a reviewed answer key.

\textit{Input/Condition:} Select a wrong answer. Wait until the reveal
finishes and any configured timeout passes; with no timeout, wait
\constref{ANSWER_RETRY_OBSERVATION_S} seconds after the reveal. Then select
a correct answer. Across attempts, cover every wrong and revealed correct
option at least once, pairing selections to avoid redundant combinations.
Also test a wrong answer followed by holding skip without a correct selection.

\textit{Output/Result:} All correct answers, and no incorrect ones, are
marked within \constref{ANSWER_REVEAL_MAX_S} second, using narration and
at least two non-colour cues. The reveal remains available for retry.
Selecting a correct answer gives encouraging feedback before continuation;
skip permits continuation without a correct answer or penalty.

\textit{Test Case Derivation:} \srsfrref{FR_RevealAnswer}{16} requires reveal and retry; \srsfrref{FR_NoPenalty}{18} permits
skipping. \srsnfrref{NFR_Colour}{10} and \srsnfrref{NFR_Responsiveness}{12} specify cues and timing, also required by SR-16
of the Hazard Analysis \cite{HazardAnalysis}.

\textit{How test will be performed:} Compare reveals with the answer key.
Record selection, reveal, retry or skip, feedback, and continuation, reusing
\testref{ST_Responsiveness} timing logs. Confirm both cues remain with
reduced motion enabled. Full colour-vision checks belong to
Section~\ref{sec:UsabilityHumanityNFRTests}; questions without wrong options
are checked in \frtestref{ST_InstructionsFeedback_06}.

\item[\frtestdef{ST_InstructionsFeedback_06}]
\textbf{Valid Questions and Build-Time Validation}

\textit{Type:} Functional, Static and Dynamic, automated validation with
manual inspection.

\textit{Requirements Covered:} \srsfrref{FR_ValidQuestions}{17}.

\textit{Initial State:} Reviewed answer keys and the build validator are
available. Prepare separate test content with one, several, or all options
correct; no correct answer; and keys naming missing options, both alone
and alongside a valid correct option.

\textit{Input/Condition:} Build with the test content, including a level
whose questions are all excluded as invalid. Then build with release content
and select each expected correct answer on a separate attempt.

\textit{Output/Result:} Every build validates its questions. Valid questions
remain playable; invalid questions are excluded and logged. Exclusion leaves
no empty question screen or blocked level. Each displayed question has at
least one correct answer, and every expected correct option is accepted.

\textit{Test Case Derivation:} \srsfrref{FR_ValidQuestions}{17} requires valid answers. Build-time
validation satisfies SR-17 of the Hazard Analysis \cite{HazardAnalysis}
and the SRS \srstableref{tab:DevelopmentPhaseThree}{Phase 3 plan}.

\textit{How test will be performed:} Check that validation runs on every
build. Compare its results and packaged question IDs with the reviewed keys;
record accepted answers and continuation after exclusions. Reuse validation
results for unchanged content, but check the installed build. Clinical
approval is covered by \testref{ST_ClinicalAccuracy}.

\item[\frtestdef{ST_InstructionsFeedback_07}]
\textbf{Skipping and Incorrect Responses without Penalties}

\textit{Type:} Functional, Dynamic, manual.

\textit{Requirements Covered:} \srsfrref{FR_NoPenalty}{18}; related
skip, completion, and reward checks under
\srsfrref{FR_DescribeSensations}{8}, \srsfrref{FR_LevelEnd}{11},
\srsfrref{FR_FirstReward}{28}, \srsfrref{FR_NoDuplicateReward}{29}.

\textit{Initial State:} For each path below, prepare an unfinished target
level with rewards from other levels where available. Record statuses,
points, coins, owned accessories, and the target's defined completion reward.

\textit{Input/Condition:} Run two paths from separate baselines: skip every
exercise and animation using the required hold; and make repeated incorrect
responses in every applicable exercise and question, then retry or skip.
Attempt to skip sensation descriptions, but let them finish. Reach the
ending, then repeat each path as a completed-level replay.

\textit{Output/Result:} All exercises and animations can be skipped except
sensation descriptions, which play in full. Skips and mistakes produce no
negative text, audio, or visual feedback and remove no points, coins, or
accessories. Both paths reach \textit{Completed} with the full defined
first-completion reward. Replay preserves completion and gives no duplicate reward.

\textit{Test Case Derivation:} \srsfrref{FR_NoPenalty}{18} prohibits penalties and blocked
completion; \srsfrref{FR_DescribeSensations}{8} protects sensation descriptions. \srsfrref{FR_LevelEnd}{11}, \srsfrref{FR_FirstReward}{28}, and \srsfrref{FR_NoDuplicateReward}{29}
define completion and rewards.

\textit{How test will be performed:} Compare rewards immediately before
and after each skip or incorrect response, then check the completion reward.
Add these checks to
runs from \frtestref{ST_InstructionsFeedback_02},
\frtestref{ST_InstructionsFeedback_05}, and
\frtestref{ST_AM_SkippedActivitiesStillComplete}. Record feedback, each
activity's skip availability, final status, and reward totals. Reuse the
sensation, clinic-variation, and speak-up message checks from
\frtestref{ST_AM_SensationBeforeSkip}, \frtestref{ST_AM_ProcedureVariation},
and \frtestref{ST_AM_ExpressConcern} on those recordings.

\item[\frtestdef{ST_InstructionsFeedback_09}]
\textbf{Repeated Touches during Transitions}

\textit{Type:} Functional, Dynamic, automated input with manual device checks.

\textit{Requirements Covered:} \srsfrref{FR_IgnoreRepeatedTouches}{19}.

\textit{Initial State:} An application screen is ready for navigation.
Logging records touch times, accepted actions, and destination display times.

\textit{Input/Condition:} Trigger each listed transition and send extra touches
on the original control, other navigation controls, and a destination
control's position. Once the destination appears, send a fresh valid input.

\textit{Output/Result:} The requested screen appears after at most one
navigation action. Extra touches neither skip screens nor execute later
on the destination. Fresh input works after the transition.

\textit{Test Case Derivation:} \srsfrref{FR_IgnoreRepeatedTouches}{19} requires additional touches to be
ignored until the next screen appears, including transitions outside levels.

\textit{How test will be performed:} Use Unity Play Mode tests through the
normal input handler. Confirm touch timing, accepted actions, and screen
order in the logs; repeat on devices with rapid taps.

\item[\frtestdef{ST_InstructionsFeedback_10}]
\textbf{Narration Protection}

\textit{Type:} Functional, Dynamic, manual.

\textit{Requirements Covered:} \srsfrref{FR_IgnoreRepeatedTouches}{19}.

\textit{Initial State:} Procedure narration is audible. An uninterrupted
recording provides its full content and ending.

\textit{Input/Condition:} Tap, hold, and repeatedly tap outside skip,
home, and restart, including other interactive elements. Repeat during
mandatory sensation descriptions.

\textit{Output/Result:} These touches neither skip nor shorten narration.

\textit{Test Case Derivation:} \srsfrref{FR_IgnoreRepeatedTouches}{19} protects procedure narration from
touches outside the three named controls.

\textit{How test will be performed:} Record touch locations and audio;
compare the content and ending with the uninterrupted recording. Home and
restart during narration are checked in
\frtestref{ST_ExitingInterruptions_01} and
\frtestref{ST_ExitingInterruptions_05}.

\item[\frtestdef{ST_InstructionsFeedback_11}]
\textbf{Skip Hold Threshold and Protected Descriptions}

\textit{Type:} Functional, Dynamic, automated input with manual device checks.

\textit{Requirements Covered:} \srsfrref{FR_IgnoreRepeatedTouches}{19},
\srsfrref{FR_DescribeSensations}{8}.

\textit{Initial State:} A skippable activity will remain active throughout
the test input, with procedure narration where present. Log touch-down,
release, and skip times. Let $T$ be \constref{SKIP_HOLD_MIN_S}.

\textit{Input/Condition:} On separate attempts, hold skip for $T-0.1$,
$T$, and $T+0.1$ seconds. Also make short presses totalling more than $T$.
Repeat during each sensation description. For each input handler, background
the app during a short press, release while away, then return and press
briefly again with the process retained.

\textit{Output/Result:} A continuous hold reaching $T$ skips a skippable
activity on the next input-processing frame. Shorter holds never skip,
and separate presses do not accumulate, including across interruptions.
Sensation descriptions cannot be skipped, even with a hold of $T$ or longer.

\textit{Test Case Derivation:} \srsfrref{FR_IgnoreRepeatedTouches}{19} defines the hold threshold and preserves
\srsfrref{FR_DescribeSensations}{8}'s exception for sensation descriptions.

\textit{How test will be performed:} Use timed Unity Play Mode input through
the normal handler. Measure actual hold durations from logs; allow the next
input-processing frame at $T$, but no early skip. Confirm short-tap and
hold behaviour on devices, and record protected descriptions finishing.

\end{frtestlist}


% ============================================================
% 5. Exiting and Interruptions
% ============================================================

\subsubsection{Exiting and Interruptions}
\label{sec:ExitingInterruptionsFRTests}

These tests verify \srsfrref{FR_ExitLevel}{20}--\srsfrref{FR_Suspend}{25}.
Use the inventory from Section~\ref{sec:InstructionsFeedbackFRTests} to
check home and restart on every screen in all three levels, including
default content and each distinct age-group variant. For \srsfrref{FR_Suspend}{25}, include
other application screens with active audio or animation.

Build an interruption matrix covering each available event type on each
distinct media or input handler, and each media sequence with at least one
applicable event. Reuse runs when the conditions match, and record untested
device events as unverified.

\begin{frtestlist}

\item[\frtestdef{ST_ExitingInterruptions_01}]
\textbf{Return to the Journey Map from Every Screen}

\textit{Type:} Functional, Dynamic, manual.

\textit{Requirements Covered:} \srsfrref{FR_ExitLevel}{20},
\srsfrref{FR_RestartLevel}{24}; one-touch exit from
\srsnfrref{NFR_Exit}{4}.

\textit{Initial State:} Setup is complete and the target level is open.
The inventory includes every screen from introduction to ending.

\textit{Input/Condition:} Select home once on each screen, including during
narration, sensation descriptions, and unfinished exercises. Re-enter the
level to reach each remaining screen.

\textit{Output/Result:} Home ends the level and returns to a usable journey
map with one touch, without confirmation or waiting for an exercise or
narration to finish.

\textit{Test Case Derivation:} \srsfrref{FR_ExitLevel}{20} defines the destination, \srsfrref{FR_RestartLevel}{24} requires
exit from every screen, and \srsnfrref{NFR_Exit}{4} sets the one-touch limit.

\textit{How test will be performed:} Record the starting screen, touch
count, and destination for each exit. Reuse these results for \srsnfrref{NFR_Exit}{4} in
Section~\ref{sec:UsabilityHumanityNFRTests}.

\item[\frtestdef{ST_ExitingInterruptions_02}]
\textbf{Level Audio and Visuals Stop on Exit}

\textit{Type:} Functional, Dynamic, manual with playback logs.

\textit{Requirements Covered:} \srsfrref{FR_StopOnExit}{21}.

\textit{Initial State:} Audio is enabled. The inventory identifies level
narration, music, sound effects, animations, and delayed or looping effects.

\textit{Input/Condition:} Exit before the ending while each media type is
active, including overlapping media. Also exit before delayed effects start
and wait on the map past their scheduled start or completion.

\textit{Output/Result:} All level media stop before the map appears and
none starts later. The map's own media may play.

\textit{Test Case Derivation:} \srsfrref{FR_StopOnExit}{21} and SR-13 of the Hazard Analysis
\cite{HazardAnalysis} require level media to stop on exit.

\textit{How test will be performed:} Reuse exits from
\frtestref{ST_ExitingInterruptions_01}. Record screen and device audio;
compare media-stop events with the first map frame in Unity logs. Check
for delayed playback and remaining active looping effects.

\item[\frtestdef{ST_ExitingInterruptions_03}]
\textbf{Early Exit Preserves Progress and Rewards}

\textit{Type:} Functional, Dynamic, manual with state inspection.

\textit{Requirements Covered:} \srsfrref{FR_ExitKeepsData}{22}.

\textit{Initial State:} Prepare a new journey and one with completed levels
and earned rewards. After starting the target level, record all level
statuses, journey points, coins, and owned accessories, including nonzero
rewards where available.

\textit{Input/Condition:} Exit an unfinished level and a completed-level
replay. Include exits after a hint, after a skip, and just before completion.

\textit{Output/Result:} All recorded values remain unchanged, with no
reward issued. The target remains \textit{In Progress} or \textit{Completed}
as appropriate; other levels are unaffected.

\textit{Test Case Derivation:} \srsfrref{FR_ExitKeepsData}{22} leaves player state unchanged on exit.
The baseline is taken after starting because \srsfrref{FR_StartLevel}{4}
can first change the target to \textit{In Progress}.

\textit{How test will be performed:} Add before-and-after state snapshots
to matching exits from \frtestref{ST_ExitingInterruptions_01}. Compare
recorded values and confirm that exit occurred before the completion event.

\item[\frtestdef{ST_ExitingInterruptions_04}]
\textbf{Start from the Beginning after Exit or Closure}

\textit{Type:} Functional, Dynamic, manual with process logs.

\textit{Requirements Covered:} \srsfrref{FR_RestartAfterExit}{23}.

\textit{Initial State:} A level is partway through play, with pending saves
complete. Record its opening screen and current step.

\textit{Input/Condition:} On separate runs, leave through home, close the
app during play, and terminate it while backgrounded. Reopen the app where
needed and start the same level. Cover unfinished levels and completed-level
replays.

\textit{Output/Result:} Each run starts at the level's beginning, without
resuming the abandoned step or resetting the journey.

\textit{Test Case Derivation:} \srsfrref{FR_RestartAfterExit}{23} requires a fresh start after exit or
closure; \srsfrref{FR_Suspend}{25} covers resuming a retained process.

\textit{How test will be performed:} Reuse exits from
\frtestref{ST_ExitingInterruptions_03}. Confirm process termination for
closure cases and compare each start with the recorded opening. Share
relaunches with \frtestref{ST_SavingRestoring_02} to check saved data.

\item[\frtestdef{ST_ExitingInterruptions_05}]
\textbf{Restart from Every Screen without Losing Progress}

\textit{Type:} Functional, Dynamic, manual with state inspection.

\textit{Requirements Covered:} \srsfrref{FR_RestartLevel}{24};
restart during narration under \srsfrref{FR_IgnoreRepeatedTouches}{19}.

\textit{Initial State:} Prepare an unfinished attempt and a completed-level
replay with earned rewards. Record saved statuses and rewards; for ending
screens, take the baseline after completion is applied.

\textit{Input/Condition:} Restart from every listed screen, including
narration, sensation descriptions, unfinished actions, hints, reveals, and
the ending. Try the first interaction and home after each restart.

\textit{Output/Result:} Restart returns to the beginning without waiting
for narration to finish, with play and home available. Statuses and rewards
are unchanged: unfinished levels remain
\textit{In Progress}, completed levels remain \textit{Completed}, and
restarting gives no reward.

\textit{Test Case Derivation:} \srsfrref{FR_RestartLevel}{24} requires restart from every screen
without losing progress; \srsfrref{FR_IgnoreRepeatedTouches}{19} keeps it available during narration.

\textit{How test will be performed:} Record the starting and returned
screens, successful interaction and exit, and before-and-after saved values.

\item[\frtestdef{ST_ExitingInterruptions_06}]
\textbf{Unexpected Inputs Do Not Block Play}

\textit{Type:} Functional, Dynamic, manual.

\textit{Requirements Covered:} \srsfrref{FR_RestartLevel}{24}.

\textit{Initial State:} For each level state, identify inputs with no defined
action. Incorrect task responses requiring hints or answer reveals are
covered by \frtestref{ST_InstructionsFeedback_02} and
\frtestref{ST_InstructionsFeedback_05}.

\textit{Input/Condition:} Apply unexpected touches and gestures once and
repeatedly, including extra simultaneous touches and cancelled drags where
undefined. Include unfinished actions. Release the inputs, then try a valid
response or permitted skip. On separate attempts, use home and restart
before completing the task.

\textit{Output/Result:} Unexpected inputs are ignored without navigation
or changes to saved progress or rewards. Valid input, home, and restart
remain usable; mandatory sensation descriptions remain unskippable.

\textit{Test Case Derivation:} \srsfrref{FR_RestartLevel}{24} requires unexpected inputs to be ignored
without blocking play; SR-4 of the Hazard Analysis \cite{HazardAnalysis}
also covers recovery from unfinished actions.

\textit{How test will be performed:} Record the state, gesture, response,
state comparison, and successful continuation or recovery. Reuse
\frtestref{ST_InstructionsFeedback_09} for transitions and
\frtestref{ST_InstructionsFeedback_10} for narration.

\item[\frtestdef{ST_ExitingInterruptions_07}]
\textbf{Pause and Resume after Backgrounding or Locking}

\textit{Type:} Functional, Dynamic, manual with playback logs.

\textit{Requirements Covered:} \srsfrref{FR_Suspend}{25}.

\textit{Initial State:} Audio is enabled and the process is retained.
Record the screen, exercise state, active media positions, level statuses,
and rewards.
Include narration, music, sound effects, and matching animations.

\textit{Input/Condition:} Background the app, then return; separately lock
and unlock the device. Wait beyond the remaining media or timer duration,
or one loop; use the same interval for input-only states. Follow the
interruption matrix, including sensation descriptions, unfinished exercises,
pending feedback/completion callbacks, and non-level media. Repeat the
pause/resume cycle on each distinct handler.

\textit{Output/Result:} All audio and animation pause without advancing.
The same screen and media positions resume with narration and visuals aligned.
Timers and callbacks do not consume paused time, skip content on return,
or complete the level during the pause. Exercises remain usable, and statuses
and rewards are unchanged.

\textit{Test Case Derivation:} \srsfrref{FR_Suspend}{25} requires continuation from the paused
point; SR-25 of the Hazard Analysis \cite{HazardAnalysis} requires media
to resume together. Start-offset timing is checked in \testref{ST_AVSync}.

\textit{How test will be performed:} Record device audio and compare media
positions at pause with those immediately before playback resumes, within
the logging precision. Compare exercise state and rewards, review resumed
content, and try valid input, plus home and restart within levels. A crash
fails the test; confirmed OS termination is checked under
\frtestref{ST_ExitingInterruptions_04}, followed by another retained-process run.

\item[\frtestdef{ST_ExitingInterruptions_08}]
\textbf{Pause and Resume after Audio or Screen Interruptions}

\textit{Type:} Functional, Dynamic, manual with playback logs.

\textit{Requirements Covered:} \srsfrref{FR_Suspend}{25}.

\textit{Initial State:} Use the states from
\frtestref{ST_ExitingInterruptions_07}, with the process retained.

\textit{Input/Condition:} Trigger available calls, alarms, or notifications
that take over audio while the game remains visible, or take over the
screen. Dismiss them and return to the game. Include sensation descriptions
and exercises awaiting input. Also test a notification that takes over
neither audio nor screen, where available.

\textit{Output/Result:} Actual takeovers meet the pause/resume criteria in
\frtestref{ST_ExitingInterruptions_07}. Notifications that take over neither
audio nor screen do not pause play.

\textit{Test Case Derivation:} \srsfrref{FR_Suspend}{25} and SR-18/SR-25 of the Hazard Analysis
\cite{HazardAnalysis} cover interruptions beyond backgrounding. Muting is
checked in \frtestref{ST_InstructionsFeedback_04}.

\textit{How test will be performed:} Confirm takeover independently of game
callbacks, then compare callback and playback logs as in
\frtestref{ST_ExitingInterruptions_07}. Use local alarms or notifications
for offline runs; record call-based and simulated checks separately from
offline device evidence.

\end{frtestlist}


% ============================================================
% 6. Saving and Restoring Progress
% ============================================================

\subsubsection{Saving and Restoring Progress}
\label{sec:SavingRestoringFRTests}

These tests verify \srsfrref{FR_PersistState}{26} and
\srsfrref{FR_CorruptData}{27} using an offline test installation.
Compare the saved level statuses and coins
(\srstableref{def_Player}{player state}), owned and equipped accessories,
and selected character. Journey points and completed-level counts are
calculated from that state, as defined by the
\srstableref{def_Scoreboard}{SRS scoreboard values}.

Keep before-and-after snapshots of the complete save data, including any
recovery copies. For \srsfrref{FR_CorruptData}{27} cases, ensure no valid fallback save can be loaded.
Keep fault-injection hooks out of release builds and repeat normal save and
restore checks on release builds. Exclude deliberately delayed or terminated
runs from normal performance and reliability measurements.

\begin{frtestlist}

\item[\frtestdef{ST_SavingRestoring_01}]
\textbf{Save Each Change before Processing the Next Input}

\textit{Type:} Functional, Dynamic, automated event checks with manual device checks.

\textit{Requirements Covered:} \srsfrref{FR_PersistState}{26}.

\textit{Initial State:} Prepare valid saves for each change: unstarted
levels for start checks, unfinished levels for completion, and sufficient
coins and an unowned item for purchases. Use the build's reward and
accessory definitions.

\textit{Input/Condition:} Separately start and first-complete each level;
purchase an accessory; equip, replace, and unequip owned items; confirm
initial and changed character selections; and confirm a journey reset
under \srsfrref{FR_ResetJourney}{42}. Include accessory rewards where defined.
Delay each save and send another valid input through the normal handler.
Use an input accepted without the delay, and hold the delay beyond other
input locks. Do not navigate away or background the app to trigger saving.

\textit{Output/Result:} Each complete change is saved before the next input
is processed. Purchases save both coin deduction and ownership; completion
saves status and rewards. Unaffected values remain unchanged, and unequipping
or resetting saves the resulting empty/default values.

\textit{Test Case Derivation:} \srsfrref{FR_PersistState}{26} requires saving each change before the
next input; checking only a later launch cannot establish this ordering.

\textit{How test will be performed:} Log the change, storage commit, and
next processed input. Delay only storage, without imposing input ordering.
Inspect saved values at commit, not when a write is merely queued. Release
the delay and try fresh input; ignored transition touches need not be queued
(\frtestref{ST_InstructionsFeedback_09}). Repeat without delays on devices,
reusing starts and completions from \frtestref{ST_AM_LevelCompletion}.

\item[\frtestdef{ST_SavingRestoring_02}]
\textbf{Restore Saved Progress and Customization across Launches}

\textit{Type:} Functional, Dynamic, manual with process and state logs.

\textit{Requirements Covered:} \srsfrref{FR_PersistState}{26}; replay and
deletion persistence under \srsfrref{FR_NoDuplicateReward}{29},
\srsfrref{FR_ResetJourney}{42}, and \srsfrref{FR_ClearAllData}{43}.

\textit{Initial State:} Prepare the mixed statuses from
\frtestref{ST_AM_JourneyMixedStatuses}, coins, owned and equipped accessories,
and a selected character. Use separate valid saves if needed, including
nonzero balances. Record all values and expected scoreboard totals.

\textit{Input/Condition:} After each save change in
\frtestref{ST_SavingRestoring_01} commits, terminate and relaunch the app.
Repeat the launch for the mixed-status save and reset state. Replay a restored
completed level to its ending. Also confirm clear-all-data and relaunch twice,
acknowledging the notice on the first launch.

\textit{Output/Result:} The latest saved values are restored and agree
with the map, scoreboard, shop, and character display. Neither launch nor
replay awards duplicate rewards. Reset and clear-all remain effective across
launches; reset keeps acknowledgement, while clear-all requires it again.

\textit{Test Case Derivation:} \srsfrref{FR_PersistState}{26} requires restoration; the SRS
\srstableref{UC_OfflineUse}{offline-use case} also covers customization.
\srsfrref{FR_NoDuplicateReward}{29} prohibits repeat rewards, and \srsfrref{FR_ResetJourney}{42}/\srsfrref{FR_ClearAllData}{43} define deletion and acknowledgement.

\textit{How test will be performed:} Verify process termination and compare
loaded and displayed values with the baseline. Include foreground and
background closure from \frtestref{ST_ExitingInterruptions_04}. Add these
checks to existing save runs. Apply \frtestref{ST_AM_ReplayCompleted} to
the restored level and the setup tests to the clear-all notice.

\item[\frtestdef{ST_SavingRestoring_03}]
\textbf{Preserve Saved Player State after a Write Failure}

\textit{Type:} Functional, Dynamic, controlled fault injection with state inspection.

\textit{Requirements Covered:} \srsfrref{FR_PersistState}{26}.

\textit{Initial State:} Copy a valid save as the baseline for each distinct
save/error-handling path in \frtestref{ST_SavingRestoring_01}. Group actions
only when both paths match, and restore the baseline before each run.

\textit{Input/Condition:} Inject storage-full or write-access failures
before writing and failures after writing begins but before commit. Cover
each path, including completion, purchase, and reset, with an accessory
reward where defined. Keep the fault active until termination to prevent
a retry or shutdown save from hiding it; remove the fault before relaunch.

\textit{Output/Result:} The error causes no crash. Saved statuses and coins
remain readable and unchanged, and relaunch restores them without partial
changes or a reset. Saving and restoration work once storage is available.

\textit{Test Case Derivation:} \srsfrref{FR_PersistState}{26} preserves previously saved player state;
SR-20 of the Hazard Analysis \cite{HazardAnalysis} requires failure handling
without crashes.

\textit{How test will be performed:} Inject errors through a storage hook
rather than filling the device. Log the failure point and compare stored
and restored values with the baseline. Record customization separately:
\srsfrref{FR_PersistState}{26} does not specify its rollback. Flag inconsistent coin/ownership results
for SRS review. After removing the fault, save a new change and relaunch.
On a separate run per error path, clear the fault without terminating;
verify usable input, including home/restart within levels, then save a new
change and relaunch to check restoration.

\item[\frtestdef{ST_SavingRestoring_04}]
\textbf{Preserve Saved Player State when a Write Is Interrupted}

\textit{Type:} Functional, Dynamic, controlled process termination with storage logs.

\textit{Requirements Covered:} \srsfrref{FR_PersistState}{26}.

\textit{Initial State:} Use the baselines and save paths from
\frtestref{ST_SavingRestoring_03}. Identify write and commit stages.

\textit{Input/Condition:} Terminate without a shutdown save before writing,
after a partial write where applicable, and immediately before commit.
Relaunch after each termination, restoring the baseline between runs.
Include a control terminated after successful commit.

\textit{Output/Result:} Interrupted writes restore the previous player
state without mixed, damaged, or reset statuses and coins. The completed-write
control restores the full new save. Relaunch does not crash.

\textit{Test Case Derivation:} \srsfrref{FR_PersistState}{26} covers interrupted writes as well as
reported failures; termination after commit must retain the new save.

\textit{How test will be performed:} Use debugger or test hooks to pause
at each storage stage, then confirm termination. Compare stored and restored
values with the baseline, or the new save for the control. Record customization
as in \frtestref{ST_SavingRestoring_03}. Unconfirmed interruption points
remain unverified.

\item[\frtestdef{ST_SavingRestoring_05}]
\textbf{Recover when Saved Data Is Missing}

\textit{Type:} Functional, Dynamic, manual with storage inspection.

\textit{Requirements Covered:} \srsfrref{FR_CorruptData}{27}.

\textit{Initial State:} Prepare a fresh installation without restored
backups, and a previously used installation with its progress and
customization data removed while closed.

\textit{Input/Condition:} Launch each installation. Capture recovery in
the load log before setup can save data, then acknowledge any required
notice without selecting a character or starting a level.

\textit{Output/Result:} Missing data is replaced without failure. Every
level is \textit{Not Started}; coins, journey points, and completed-level
count are zero. No accessories are owned/equipped and no character is selected.
All levels remain available using the default character.

\textit{Test Case Derivation:} \srsfrref{FR_CorruptData}{27} specifies new-installation progress
and customization; \srsfrref{FR_ChooseCharacter}{31} provides a default character.

\textit{How test will be performed:} Confirm missing data before launch
and inspect the replacement. Reuse \frtestref{ST_AM_JourneyFresh} for map
checks, then start a level, wait for saving, and relaunch to confirm new
progress persists. Retain snapshots and load logs.

\item[\frtestdef{ST_SavingRestoring_06}]
\textbf{Recover when Saved Data Cannot Be Read}

\textit{Type:} Functional, Dynamic, controlled fault injection with state inspection.

\textit{Requirements Covered:} \srsfrref{FR_CorruptData}{27}.

\textit{Initial State:} Valid saved progress and customization exist and
have been copied. The app is closed; storage remains writable.

\textit{Input/Condition:} Launch with a hook causing an access or I/O error
when reading the existing save, while allowing replacement data to be written.

\textit{Output/Result:} The app replaces unreadable data with the defaults
in \frtestref{ST_SavingRestoring_05} without failure or partial recovery.
Later launches use the replacement rather than restoring the old data.

\textit{Test Case Derivation:} \srsfrref{FR_CorruptData}{27} separately requires recovery from
existing data that cannot be read.

\textit{How test will be performed:} Confirm the save exists and the read
error reaches the load path. Inspect replacement data, remove the fault,
and relaunch. Repeat the continued-play/save check in
\frtestref{ST_SavingRestoring_05}, distinguishing injected errors from
reproduced device failures.

\item[\frtestdef{ST_SavingRestoring_07}]
\textbf{Recover when Saved Data Is Invalid}

\textit{Type:} Functional, Dynamic, prepared invalid data with device checks.

\textit{Requirements Covered:} \srsfrref{FR_CorruptData}{27}.

\textit{Initial State:} Copy a valid save with progress and customization.
Identify required fields and allowed empty values in the storage format.
Change one condition per invalid copy, leaving other fields valid.

\textit{Input/Condition:} Launch separately with empty, malformed, or
truncated data; missing required fields; wrong types; negative coins;
unknown statuses or character/accessory IDs; missing, unknown, or duplicate
level records where representable; unowned equipped items; two equipped
items in one category; and a reward-claimed flag inconsistent with completion
status if stored.
Include an earlier-format save only if current validation and any supported
migration reject it.

\textit{Output/Result:} Validation rejects invalid data before gameplay or
progress display. The whole progress/customization state is replaced with
the defaults in \frtestref{ST_SavingRestoring_05}, without failure or
retained old values. Replacement data remains usable across launches.

\textit{Test Case Derivation:} \srsfrref{FR_CorruptData}{27} requires replacement; SR-20 of the
Hazard Analysis \cite{HazardAnalysis} requires validation before use.
Valid values follow the \srstableref{def_Player}{player-state model},
\srstableref{tab:JourneyDataDictionary}{data dictionary}, and
\srsfrref{FR_Equip}{36}.

\textit{How test will be performed:} Restore the baseline for each case.
Record the mutation and rejection in load logs; compare replacement data
and displays with the defaults, then repeat the continued-play/save check
in \frtestref{ST_SavingRestoring_05}. Use valid populated and new-player saves
as controls, including unset optional fields, empty inventories, and completed
levels with zero coins after spending where reachable. Valid controls must
load without recovery.

\end{frtestlist}


% ============================================================
% 7. Rewards and Scoreboard
% ============================================================

\subsubsection{Rewards and Scoreboard}
\label{sec:RewardsScoreboardFRTests}

These tests verify the level rewards and the personal scoreboard in
\srsfrref{FR_FirstReward}{28}--\srsfrref{FR_Scoreboard}{30} of the SRS
\cite{SRS}. They use the test levels and saved states in
Appendix~\ref{sec:RewardFixtures}, so the expected values do not change when
the real rewards are tuned. Each test is an automated Unity Play Mode test
that performs the actions and then checks the saved game state.

\begin{frtestlist}

\item[\frtestdef{ST_RewardsScoreboard_01}]\textbf{Coin Reward on First Completion}

\textit{Type:} Functional, Dynamic, automatic.

\textit{Requirements Covered:} \srsfrref{FR_FirstReward}{28}.

\textit{Initial State:} The saved state is \constref{SAVE_NEW}.

\textit{Input/Condition:} Start \constref{L1}, skip its exercise, and answer
its question incorrectly before selecting the correct answer. Reach the end
of the level.

\textit{Output/Result:} The reward screen shows 10 journey points and 5 coins
earned. The level is Completed, the coin balance is 5, and the total journey
points are 10.

\textit{Test Case Derivation:} \srsfrref{FR_FirstReward}{28} gives the level's
journey points and coins when it is finished for the first time. The reward
depends on finishing the level, not on the answers, so skipping and a wrong
answer do not reduce it.

\textit{How test will be performed:} The Play Mode test plays the level, reads
the values on the reward screen, and then checks the saved level status, coin
balance, and total journey points.

\item[\frtestdef{ST_RewardsScoreboard_02}]\textbf{Accessory Reward on First Completion}

\textit{Type:} Functional, Dynamic, automatic.

\textit{Requirements Covered:} \srsfrref{FR_FirstReward}{28}.

\textit{Initial State:} The saved state is \constref{SAVE_NEW}.

\textit{Input/Condition:} Complete \constref{L3}. Do not equip the accessory
that is offered.

\textit{Output/Result:} The reward screen shows 40 journey points, 15 coins,
and the Sun hat, with an option to equip it. The Sun hat is owned but not
equipped. The coin balance is 15.

\textit{Test Case Derivation:} \srsfrref{FR_FirstReward}{28} includes the
level's accessory in the reward when its Reward Definition has one.
\srsitemref{UC_CompleteGame}{UC9} offers to equip the accessory but does not
equip it.

\textit{How test will be performed:} The Play Mode test completes the level,
reads the reward screen, and then checks the saved coin balance and the
owned and equipped accessories.

\item[\frtestdef{ST_RewardsScoreboard_03}]\textbf{No Duplicate Reward on Replay}

\textit{Type:} Functional, Dynamic, automatic.

\textit{Requirements Covered:} \srsfrref{FR_NoDuplicateReward}{29}.

\textit{Initial State:} The saved state is \constref{SAVE_MIXED}.

\textit{Input/Condition:} Replay \constref{L3} to its end
\constref{REPLAY_TRIALS_MIN} times. Restart the game session before the last
replay.

\textit{Output/Result:} Each replay ends with a congratulation and no points,
coins, or accessory. The total journey points stay at 50, the coin balance
stays at 3, and the owned accessories do not change.

\textit{Test Case Derivation:} \srsfrref{FR_NoDuplicateReward}{29} forbids a
second reward for a level that has already been rewarded. The coins from
\constref{L3} were already spent in \constref{SAVE_MIXED}, so a duplicate
reward would raise the balance to 18.

\textit{How test will be performed:} The Play Mode test replays the level and
checks the end-of-level screen and saved game state after each replay. It
restarts the session by reloading the game from the save file.

\item[\frtestdef{ST_RewardsScoreboard_04}]\textbf{Scoreboard for a New Player}

\textit{Type:} Functional, Dynamic, automatic.

\textit{Requirements Covered:} \srsfrref{FR_Scoreboard}{30}.

\textit{Initial State:} The saved state is \constref{SAVE_NEW}.

\textit{Input/Condition:} Open the scoreboard from the journey map.

\textit{Output/Result:} Each test level is shown as Not Started with 0 points.
The scoreboard shows 0 total journey points, 0 of 3 levels completed, and
0 coins. No nickname is shown, and the wording is not personalized.

\textit{Test Case Derivation:} A new player has no completed levels and no
coins, so every scoreboard value is 0. \srsfrref{FR_Nickname}{40} requires
wording that is not personalized when no nickname is set.

\textit{How test will be performed:} The Play Mode test opens the scoreboard
and checks each value shown against the expected result.

\item[\frtestdef{ST_RewardsScoreboard_05}]\textbf{Scoreboard with Progress}

\textit{Type:} Functional, Dynamic, automatic.

\textit{Requirements Covered:} \srsfrref{FR_Scoreboard}{30}.

\textit{Initial State:} The saved state is \constref{SAVE_MIXED}.

\textit{Input/Condition:} Open the scoreboard, return to the journey map, and
open it again.

\textit{Output/Result:} \constref{L1} is Completed with 10 points,
\constref{L2} is In Progress with 0 points, and \constref{L3} is Completed
with 40 points. The scoreboard shows 50 total journey points, 2 of 3 levels
completed, 3 coins, the nickname Sam, and Test character B. It shows no rank
or results from other players. Viewing it does not change the saved game
state.

\textit{Test Case Derivation:} \srsfrref{FR_Scoreboard}{30} counts points only
for completed levels, so the total is $10 + 40 = 50$. The coin balance is
stored separately and is 3 after the Blue scarf was bought. The scoreboard
must not compare the child with other users.

\textit{How test will be performed:} The Play Mode test opens the scoreboard
twice and checks each value shown. It compares the saved game state before
and after viewing the scoreboard.

\end{frtestlist}


% ============================================================
% 8. Characters and Accessories
% ============================================================

\subsubsection{Characters and Accessories}
\label{sec:CharactersAccessoriesFRTests}

These tests verify the character selection and customization
features defined in the SRS. They cover how players interact
with available characters and accessories and how their
selections affect the game.

\begin{frtestlist}

\item[\frtestdef{ST_CharactersAccessories_01}]
\textbf{[Test Case Name]}

\textit{Type:} Functional, Dynamic, manual.

\textit{Requirements Covered:}
[Insert SRS functional requirement reference(s).]

\textit{Initial State:}
[Describe the available characters, accessories, and
customization state.]

\textit{Input/Condition:}
[Describe the character or accessory interaction being tested.]

\textit{Output/Result:}
[Describe the expected customization behaviour.]

\textit{Test Case Derivation:}
[Explain how the test relates to the specified requirement.]

\textit{How test will be performed:}
[Describe the steps and evidence used to verify the outcome.]

% Add additional test cases as needed.

\end{frtestlist}


% ============================================================
% 9. Caregiver Features
% ============================================================

\subsubsection{Caregiver Features}
\label{sec:CaregiverFeaturesFRTests}

These tests verify the caregiver-facing features defined
in the SRS. They cover access to caregiver information,
application preferences, and other caregiver controls,
ensuring that the available functionality behaves as
specified.

\begin{frtestlist}

\item[\frtestdef{ST_CaregiverFeatures_01}]
\textbf{[Test Case Name]}

\textit{Type:} Functional, Dynamic, manual.

\textit{Requirements Covered:}
[Insert SRS functional requirement reference(s).]

\textit{Initial State:}
[Describe the caregiver settings and application state.]

\textit{Input/Condition:}
[Describe the caregiver action or feature being tested.]

\textit{Output/Result:}
[Describe the expected caregiver feature behaviour.]

\textit{Test Case Derivation:}
[Explain how the test verifies the corresponding SRS requirement.]

\textit{How test will be performed:}
[Describe how the feature will be tested and results recorded.]

% Add additional test cases as needed.

\end{frtestlist}

\subsection{Tests for Nonfunctional Requirements}\label{sec:NFRSystemTests}

The tests in this section verify the non-functional requirements (NFRs)
defined in the SRS. These requirements describe how well the game
should operate, rather than the specific features it provides.
The tests are organized according to the NFR categories in the
SRS, with each test identifying the requirement it covers, the conditions
for testing, the expected result, and how the test will be performed. Test
parameters are listed in Section~\ref{sec:VnVParameters}.

For the Capstone course, the team will test on the limited physical devices
available and use Unity tools to simulate additional device configurations
where appropriate. The devices or simulations used will be recorded when
the tests are performed. Any additional changes to the testing scope for
the Capstone course will be explained in a \textit{Remark} under the
relevant test.


% ============================================================
% 8.1 Look and Feel Requirements
% ============================================================

\subsubsection{Look and Feel Requirements}
\label{sec:LookAndFeelNFRTests}

These tests verify the control layout and the content safety requirements in
\srsnfrref{NFR_Layout}{1}--\srsnfrref{NFR_ChildSafe}{2} of the SRS
\cite{SRS}. The survey in Appendix~\ref{sec:UsabilitySurvey} adds participant
feedback on both requirements.

\begin{testlist}

\item[\testdef{ST_Layout}]\textbf{Consistent Control Positions}

\textit{Type:} Non-functional, Dynamic, automatic.

\textit{Requirements Covered:} \srsnfrref{NFR_Layout}{1}.

\textit{Initial State:} The latest build is open in the Unity editor. The designated position of each shared control, such as the exit, hint, and settings buttons, is recorded in the design.

\textit{Input/Condition:} Visit every screen of every level, the journey map, the scoreboard, and the character screen at each test screen size.

\textit{Output/Result:} On every screen where a shared control appears, its centre is within \constref{CONTROL_POSITION_TOLERANCE_MM} millimetres of its designated position.

\textit{How test will be performed:} A Play Mode test opens each screen and records the on-screen position of each shared control. It converts the distance from the designated position to millimetres using the screen density of each test screen size. A team member reviews any screen that fails and checks the results on one physical device.

\item[\testdef{ST_ChildSafe}]\textbf{Child-Safe Content Review}

\textit{Type:} Non-functional, Static content inspection.

\textit{Requirements Covered:} \srsnfrref{NFR_ChildSafe}{2}.

\textit{Review procedure:} One team member will prepare a list of every image, animation, narration clip, and sound in the release build, with a recording of each level. A child-life specialist will review each item and answer two questions. Is the item frightening for a child aged 3 to 7? Does it show pain or distress beyond the sensations in the approved script? A second team member will record each finding in GitHub Issues. Changed content will be sent back for review.

\textit{Expected result and evidence:} Every item on the list has been reviewed, and no finding is unresolved at release. The team will keep the content list, the completed checklist, and the related GitHub Issues.

\textit{Remark:} If a child-life specialist is not available during the capstone, the project supervisor will review the content instead. This would not fully verify the requirement, since the SRS names a child-life specialist as the reviewer.

\end{testlist}

% ============================================================
% 8.2 Usability and Humanity Requirements
% ============================================================

\subsubsection{Usability and Humanity Requirements}
\label{sec:UsabilityHumanityNFRTests}

These tests verify the usability and humanity requirements in
\srsnfrref{NFR_TouchTarget}{3}--\srsnfrref{NFR_Colour}{10} of the SRS
\cite{SRS}. Children and caregivers will not take part in testing during the
capstone (Section~\ref{sec:Objectives}). The tests that need users will use
classmates and clinical staff instead, as explained in the remark under each
test. The survey in Appendix~\ref{sec:UsabilitySurvey} adds participant feedback
on these requirements.

\begin{testlist}

\item[\testdef{ST_TouchTarget}]\textbf{Touch Target Size}

\textit{Type:} Non-functional, Dynamic, automatic with manual check.

\textit{Requirements Covered:} \srsnfrref{NFR_TouchTarget}{3}.

\textit{Initial State:} The latest build is open in the Unity editor and installed on the test devices.

\textit{Input/Condition:} Visit every screen of the game at each test screen size.

\textit{Output/Result:} Every interactive element has a touch target at least \constref{TOUCH_TARGET_MIN} millimetres wide and high.

\textit{How test will be performed:} A Play Mode test opens each screen, finds every interactive element, and converts the size of its touch area to millimetres using the screen density. A team member then measures the smallest elements on each physical test device with a ruler or a screenshot of known scale.

\item[\testdef{ST_Exit}]\textbf{One-Touch Exit from Every Screen}

\textit{Type:} Non-functional, Dynamic, automatic.

\textit{Requirements Covered:} \srsnfrref{NFR_Exit}{4}.

\textit{Initial State:} The latest build is open in the Unity editor.

\textit{Input/Condition:} Open each level, and on every screen of the level, touch the exit button once.

\textit{Output/Result:} On every screen of every level, one touch shows the journey map.

\textit{How test will be performed:} A Play Mode test steps through each level screen by screen. On each screen, it touches the exit button once and checks that the journey map is shown. It then reopens the level and continues from the next screen.

\item[\testdef{ST_Usability}]\textbf{Core Tasks Without Instructions}

\textit{Type:} Non-functional, Dynamic, manual.

\textit{Requirements Covered:} \srsnfrref{NFR_Usability}{5}.

\textit{Initial State:} The latest build is installed on a test device. The participant has not used the game before and has not been shown how it works.

\textit{Input/Condition:} At least \constref{USABILITY_PARTICIPANTS_MIN} participants each complete the three core tasks. They start a named level from the journey map, play a level to its end, and leave a level early. Each participant acts as one of the child personas in the SRS.

\textit{Output/Result:} At least \constref{UNAIDED_RATE_MIN}\% of participants complete each task without verbal help. The median time for each task is at most \constref{CORE_TASK_MAX_MIN} minutes.

\textit{How test will be performed:} A team member gives each task in one sentence and then does not help. They record whether the task was completed without help and how long it took. They also note where the participant hesitated or made a wrong touch. The participant then completes the survey in Appendix~\ref{sec:UsabilitySurvey}.

\textit{Remark:} The participants will be classmates rather than children, because the capstone does not include testing with children. This finds usability problems but does not verify the requirement for the intended users.

\item[\testdef{ST_ClinicianUse}]\textbf{Clinician Finds a Named Level}

\textit{Type:} Non-functional, Dynamic, manual.

\textit{Requirements Covered:} \srsnfrref{NFR_ClinicianUse}{6}.

\textit{Initial State:} The latest build is installed on a test device and is open on the home screen. The participant is a clinician or clinical staff member who has not used the game before.

\textit{Input/Condition:} A team member trains the participant for at most \constref{CLINICIAN_TRAINING_MAX_MIN} minutes. The participant is then handed the device and asked to open the introduction of a named level.

\textit{Output/Result:} Every participant opens the introduction within \constref{CLINICIAN_FIND_MAX_MIN} minute of picking up the device.

\textit{How test will be performed:} A team member times the training and the task, and names a different level for each participant. They record the time taken and any wrong touches. The participant then completes the survey in Appendix~\ref{sec:UsabilitySurvey}.

\textit{Remark:} The number of participants depends on the clinical staff the project supervisor can arrange. If only the supervisor takes part, the result will be recorded as preliminary.

\item[\testdef{ST_NoPressure}]\textbf{No Pressure to Keep Playing}

\textit{Type:} Non-functional, Static design and build inspection.

\textit{Requirements Covered:} \srsnfrref{NFR_NoPressure}{7}.

\textit{Review procedure:} Two team members will separately go through every screen of the release build with a checklist of the patterns in the requirement. These are rewards for consecutive days, rewards or offers that expire, countdown timers outside an exercise, randomized rewards, and notifications. They will also check that the Android manifest and iOS project settings do not request permission to send notifications. Any pattern found will be recorded in GitHub Issues and checked again after it is fixed.

\textit{Expected result and evidence:} The checklist finds none of the patterns, and neither build requests notification permission. The team will keep the completed checklists and the build settings checked.

\item[\testdef{ST_Language}]\textbf{Sentence Length and Reading Level}

\textit{Type:} Non-functional, Static text inspection, automatic with manual review.

\textit{Requirements Covered:} \srsnfrref{NFR_Language}{8}.

\textit{Review procedure:} A script will count the words in every narrated sentence of the narration scripts. It will also compute the Flesch--Kincaid grade of each text in the caregiver area. A child-life specialist will then read the narration scripts and list any word a child aged 3 to 7 may not know. Each sentence or word that fails will be rewritten and checked again.

\textit{Expected result and evidence:} No narrated sentence has more than \constref{SENTENCE_MAX_WORDS} words, no unfamiliar word is unresolved, and every caregiver-area text has a grade of at most \constref{ADULT_TEXT_MAX_GRADE}. The team will keep the script output and the specialist's word list.

\textit{Remark:} If a child-life specialist is not available, the project supervisor will review the words instead, as in \testref{ST_ChildSafe}.

\item[\testdef{ST_Accessibility}]\textbf{WCAG 2.2 Level AA Audit}

\textit{Type:} Non-functional, Static accessibility audit with manual in-app checks.

\textit{Requirements Covered:} \srsnfrref{NFR_Accessibility}{9}.

\textit{Review procedure:} One team member will list the WCAG 2.2 Level AA success criteria that apply to a native mobile game, with a reason for each criterion left out. They will check each criterion on the test devices, using a contrast checker for text and controls and the screen reader for labels. A second team member will repeat the audit separately. Each failure will be recorded in GitHub Issues and checked again after it is fixed.

\textit{Expected result and evidence:} The audit finds no failure of an applicable criterion. The team will keep the list of criteria, both completed audits, and the related GitHub Issues.

\item[\testdef{ST_Colour}]\textbf{Screens Without Colour}

\textit{Type:} Non-functional, Dynamic, manual with simulation.

\textit{Requirements Covered:} \srsnfrref{NFR_Colour}{10}.

\textit{Initial State:} The latest build is open in the Unity editor with a colour vision deficiency filter.

\textit{Input/Condition:} View every screen under protanopia, deuteranopia, and tritanopia simulation. Include the journey map with all three level statuses and a revealed correct answer in each level.

\textit{Output/Result:} Every state and choice can be told apart by shape, position, or text as well as colour. Every revealed answer has at least two cues that do not rely on colour, and both are visible in all three simulations.

\textit{How test will be performed:} A team member takes a screenshot of each screen under each simulation. A second team member reviews the screenshots and records any state or choice that can only be told apart by colour.

\end{testlist}


\subsubsection{Performance Requirements}\label{sec:PerformanceNFRTests}

These tests verify the performance requirements defined
in the SRS. They cover startup time, responsiveness,
animation and audio behaviour, reliability, and resource
usage during normal gameplay.


\begin{testlist}

\item[\testdef{ST_Startup}]\textbf{Startup Time}

\textit{Type:} Non-functional, Dynamic, manual.

\textit{Requirements Covered:} \srsnfrref{NFR_Startup}{11}.

\textit{Initial State:} The latest build is installed on the test devices, and the app is closed.

\textit{Input/Condition:} Open the app from its icon at least \constref{STARTUP_TRIALS_MIN} times on each test device. Close the app fully between launches.

\textit{Output/Result:} The median time to reach the usable home screen is at most \constref{STARTUP_MAX_S} seconds on each device tested.

\textit{How test will be performed:} Record each launch and measure the time from tapping the icon to when the home screen is ready. Calculate the median for each device.

\item[\testdef{ST_Responsiveness}]\textbf{Touch Response}

\textit{Type:} Non-functional, Dynamic, manual with timing logs.

\textit{Requirements Covered:} \srsnfrref{NFR_Responsiveness}{12}.

\textit{Initial State:} The app is open on the test devices.

\textit{Input/Condition:} Make at least \constref{TOUCH_TRIALS_MIN} valid touches across the menus and games. Also test each type of incorrect answer and time the correct-answer reveal.

\textit{Output/Result:} On each device tested, at least \constref{RESPONSE_RATE_MIN}\% of touches show a response within \constref{RESPONSE_MAX_MS} milliseconds. Every correct-answer reveal takes at most \constref{ANSWER_REVEAL_MAX_S} second.

\textit{How test will be performed:} Use Unity timestamps to record each touch event and its first visible response. Calculate the percentage that meets the limit, and record the time taken to reveal the correct answer after an incorrect response.

\item[\testdef{ST_Smoothness}]\textbf{Animation Frame Rate}

\textit{Type:} Non-functional, Dynamic, automatic.

\textit{Requirements Covered:} \srsnfrref{NFR_Smoothness}{13}.

\textit{Initial State:} The latest build is installed on the test devices.

\textit{Input/Condition:} Play through each procedure walkthrough and practice game while recording frame times.

\textit{Output/Result:} On each device tested, the game runs at or above \constref{FRAME_RATE_MIN} frames per second for at least \constref{FRAME_RATE_COMPLIANCE_MIN}\% of the measured time.

\textit{How test will be performed:} Use the Unity Profiler to collect frame times throughout each walkthrough and game. Calculate the percentage of measured time that meets the required frame rate.

\item[\testdef{ST_AVSync}]\textbf{Narration and Animation Timing}

\textit{Type:} Non-functional, Dynamic, automatic with manual review.

\textit{Requirements Covered:} \srsnfrref{NFR_AVSync}{14}.

\textit{Initial State:} The latest narration and animations are included in the build.

\textit{Input/Condition:} Play each narrated step in every procedure walkthrough on the test devices.

\textit{Output/Result:} The narration starts within \constref{AV_SYNC_MAX_MS} milliseconds of its matching animation for every step.

\textit{How test will be performed:} Add temporary Unity logs at the start of each narration clip and its animation. Compare the timestamps and review a recording of any step that appears out of sync.

\item[\testdef{ST_AudioSafety}]\textbf{Audio Levels}

\textit{Type:} Non-functional, Static audio inspection.

\textit{Requirements Covered:} \srsnfrref{NFR_AudioSafety}{15}.

\textit{Review procedure:} One team member will export the final narration, sound effects, and music with the game's audio mixer settings applied, then compare their peak levels using Audacity or a similar audio editor. A second team member will check the measurements and mixer settings. Any sound above the narration peak will be adjusted and checked again.

\textit{Expected result and evidence:} No sound effect or music has a peak level above narration. The team will save the measurement results, mixer settings, and any issues corrected.

\item[\testdef{ST_Flashing}]\textbf{Flashing Content}

\textit{Type:} Non-functional, Dynamic, automated analysis with manual review.

\textit{Requirements Covered:} \srsnfrref{NFR_Flashing}{16}.

\textit{Initial State:} All animations, transitions, and reward effects are included in the build.

\textit{Input/Condition:} Record or export every animation and visual effect, including repeated or looping effects.

\textit{Output/Result:} No sequence contains more than \constref{FLASH_MAX_PER_S} flashes in any one-second period.

\textit{How test will be performed:} A team member will run the recordings through a photosensitive-epilepsy flash-analysis tool. A second team member will review any flagged sections. Save the results and rerun the analysis after fixing any failing effects.

\item[\testdef{ST_Reliability}]\textbf{Crash-Free Sessions}

\textit{Type:} Non-functional, Dynamic, manual.

\textit{Requirements Covered:} \srsnfrref{NFR_Reliability}{17}.

\textit{Initial State:} The current build is installed on the test devices and crash logs are available.

\textit{Input/Condition:} Complete at least \constref{TEST_SESSIONS_MIN} sessions across the test devices. Include normal gameplay, restarting levels, saving, closing and reopening, and suspend/resume.

\textit{Output/Result:} At least \constref{CRASH_FREE_RATE_MIN}\% of sessions finish without a crash or freeze.

\textit{How test will be performed:} Record the result of each session. Check crash logs and count the app as frozen if it does not respond to a normal action for \constref{FREEZE_OBSERVATION_S} seconds. Calculate the percentage of sessions completed without a crash or freeze.

\item[\testdef{ST_DeviceHeat}]\textbf{CPU Usage During Gameplay}

\textit{Type:} Non-functional, Dynamic, automatic.

\textit{Requirements Covered:} \srsnfrref{NFR_DeviceHeat}{18}.

\textit{Initial State:} The app is installed on the test devices.

\textit{Input/Condition:} Play continuously for at least \constref{HEAT_TEST_SESSION_MIN} minutes on each test device, moving between the different activities.

\textit{Output/Result:} Measured CPU utilization stays at or below \constref{CPU_UTILIZATION_MAX_PERCENT}\% throughout the session on each device tested.

\textit{How test will be performed:} Use the Unity Profiler and available device monitoring tools to record CPU utilization during gameplay. Record the highest measured value and any thermal warning.

\end{testlist}

\subsubsection{Operational and Environmental Requirements}\label{sec:OperationalNFRTests}

These tests verify the operational and environmental
requirements defined in the SRS. They check that the
game can be installed and used offline, meets the storage
limit, and includes clinically approved content.


\begin{testlist}

\item[\testdef{ST_Offline}]\textbf{Offline Functionality}

\textit{Type:} Non-functional, Dynamic, manual.

\textit{Requirements Covered:} \srsnfrref{NFR_Offline}{19}.

\textit{Initial State:} The game is installed with all content, but has not been opened for the first time.

\textit{Input/Condition:} Turn on airplane mode and disable Wi-Fi. Run all functional tests without connecting to the internet.

\textit{Output/Result:} Every functional test passes without requiring a network connection.

\textit{How test will be performed:} Run the full set of functional tests on the test devices and record the results. Record any missing content, error message, or feature that cannot be used offline.

\item[\testdef{ST_Installability}]\textbf{Installation and First Launch}

\textit{Type:} Non-functional, Dynamic, manual.

\textit{Requirements Covered:} \srsnfrref{NFR_Installability}{20}.

\textit{Initial State:} The app is not installed on the test device.

\textit{Input/Condition:} Install the latest release build and open the app for the first time.

\textit{Output/Result:} The home screen can be reached without creating an account, granting permissions, or downloading more content.

\textit{How test will be performed:} Install the build on each test device, record the first launch, and note any prompt or extra step required before reaching the home screen.

\textit{Remark:} For this capstone, the build will be installed directly for testing rather than through a public app-store listing. The store-installation part of the SRS requirement will remain unverified unless the app is published.

\item[\testdef{ST_InstallSize}]\textbf{Installed Application Size}

\textit{Type:} Non-functional, Static build inspection.

\textit{Requirements Covered:} \srsnfrref{NFR_InstallSize}{21}.

\textit{Review procedure:} Once the iOS and Android release builds are ready, the team member responsible for each build will check the installed application size using the platform's storage information. A second team member will check the reported numbers against \constref{INSTALL_SIZE_MAX_MB} MB. The build version, measured size, and source of the measurement will be recorded.

\textit{Expected result and evidence:} The installed size of each build is no more than \constref{INSTALL_SIZE_MAX_MB} MB. The team will keep screenshots or build reports showing the size.

\textit{Remark:} The SRS fit criterion uses the size reported by each app store. For the capstone, the team will use locally measured installed size unless a store release is available, and record that difference in the results.

\item[\testdef{ST_ClinicalAccuracy}]\textbf{Clinical Review of Game Content}

\textit{Type:} Non-functional, Manual document review.

\textit{Requirements Covered:} \srsnfrref{NFR_ClinicalAccuracy}{22}.

\textit{Review procedure:} One team member will prepare the latest game walkthroughs, narration, and procedure guide, with the version of each item recorded. A paediatric ophthalmologist will review the content against the procedures used at the clinic. Another team member will record the feedback and any required changes. Updated content will be sent back for approval.

\textit{Expected result and evidence:} The ophthalmologist approves every level and procedure-guide entry for the version being released, with no unresolved concerns. The team will keep the content checklist and written approval.

\end{testlist}

\subsubsection{Maintainability and Support Requirements}\label{sec:MaintainabilityNFRTests}

These tests verify the maintainability and support
requirements defined in the SRS. They cover the effort
needed to make likely changes, compatibility with supported
devices, and testing after third-party dependency updates.


\begin{testlist}

\item[\testdef{ST_Maintainability}]\textbf{Effort to Make a Likely Change}

\textit{Type:} Non-functional, Manual estimation review.

\textit{Requirements Covered:} \srsnfrref{NFR_Maintainability}{23}.

\textit{Review procedure:} Once the main game is implemented, each team member will estimate how many hours it would take to implement and test each likely change listed in the SRS, excluding clinical review. The team will discuss the estimates and agree on a final number for each change. These will be compared with the original development effort.

\textit{Expected result and evidence:} Each estimated change effort is less than \constref{MAINT_FRACTION} of the original development effort. The team will keep the individual estimates, agreed estimates, and comparisons.

\textit{Remark:} For the scope of the capstone, the team will estimate the effort rather than carry out every likely change. This gives an indication of maintainability but does not fully verify the SRS criterion based on recorded effort.

\item[\testdef{ST_Portability}]\textbf{iOS and Android Compatibility}

\textit{Type:} Non-functional, Dynamic, manual.

\textit{Requirements Covered:} \srsnfrref{NFR_Portability}{24}.

\textit{Initial State:} The latest build is ready and the test configurations have been recorded.

\textit{Input/Condition:} Run all functional tests on the supported test configurations and check the game's layouts and controls at screen sizes from \constref{SCREEN_MIN_IN} to \constref{SCREEN_MAX_IN} inches.

\textit{Output/Result:} All functional tests pass on the configurations tested, and the layouts and controls remain usable at the tested screen sizes.

\textit{How test will be performed:} Run the tests and record the device or configuration, OS version, screen size, and results. List any supported iOS or Android versions or screen sizes that could not be checked.

\item[\testdef{ST_DependencyTesting}]\textbf{Third-Party Dependency Checks}

\textit{Type:} Non-functional, Dynamic regression testing with dependency inspection.

\textit{Requirements Covered:} \srsnfrref{NFR_DependencyTesting}{25}.

\textit{Initial State:} The Unity project, dependency files, and release build are available.

\textit{Input/Condition:} Check the Unity version and all third-party package and plugin versions. Run all functional tests on the release build after each dependency change.

\textit{Output/Result:} Every dependency has an exact version, and all functional tests pass after each dependency change with no unresolved issue before release.

\textit{How test will be performed:} One team member will prepare the dependency list from the Unity Package Manager and package files. A second team member will check the versions against the project. After a version change, the team will run the functional tests and keep the results with the dependency list.

\end{testlist}

\subsubsection{Security Requirements}\label{sec:SecurityNFRTests}

These tests verify the security requirements defined
in the SRS. They check how the game stores and protects
local data, whether it collects only the data it needs,
and whether its dependencies have known security issues.


\begin{testlist}

\item[\testdef{ST_DataProtection}]\textbf{Private Local Storage}

\textit{Type:} Non-functional, Dynamic, manual.

\textit{Requirements Covered:} \srsnfrref{NFR_DataProtection}{26}.

\textit{Initial State:} A test nickname, settings, and game progress have been saved on an unmodified test device.

\textit{Input/Condition:} Try accessing the saved files from the device file manager, a different app, and normal sharing options.

\textit{Output/Result:} The saved data cannot be read by another app or through device file-sharing features.

\textit{How test will be performed:} Confirm where Unity writes the saved data and inspect the app's storage and sharing settings. Record the access attempts and check that the game can still load its own saves.

\item[\testdef{ST_Integrity}]\textbf{Release Signing and Distribution}

\textit{Type:} Non-functional, Static release inspection.

\textit{Requirements Covered:} \srsnfrref{NFR_Integrity}{27}.

\textit{Review procedure:} If the app is prepared for public release, the team members responsible for the iOS and Android builds will check the signatures using Xcode and Android build tools. Another team member will review the signing results and confirm that the planned public distribution channels are the Apple App Store and Google Play Store. The team will save the signing reports and store-release records, if available.

\textit{Expected result and evidence:} Both release builds have valid signatures, and public distribution uses only the official stores.

\textit{Remark:} This test will likely be omitted during the capstone because publishing the app to the Apple App Store and Google Play Store is difficult within the course timeline. The team may check signing on test builds, but that would not verify the full release and distribution requirement.

\item[\testdef{ST_Privacy}]\textbf{Network Activity}

\textit{Type:} Non-functional, Dynamic, manual with network capture.

\textit{Requirements Covered:} \srsnfrref{NFR_Privacy}{28}.

\textit{Initial State:} The game is installed on a test device connected to a monitored Wi-Fi network.

\textit{Input/Condition:} Capture network traffic while launching the game, playing each level, changing settings, earning rewards, and closing and reopening the game.

\textit{Output/Result:} The game makes no outgoing network requests.

\textit{How test will be performed:} Capture traffic using a network monitor or test proxy and look for requests associated with the app. Also check the Unity project for enabled analytics or network services. Keep the capture and record any unexplained request.

\item[\testdef{ST_Minimization}]\textbf{Stored Data Fields}

\textit{Type:} Non-functional, Dynamic save-file inspection.

\textit{Requirements Covered:} \srsnfrref{NFR_Minimization}{29}.

\textit{Initial State:} The app is installed with a fresh save file.

\textit{Input/Condition:} Set a test nickname, complete activities, earn and spend coins, and change the available settings. Then inspect the stored files.

\textit{Output/Result:} The app stores only fields belonging to Local Preferences, Character Inventory, Game Progress, and Journey Progress in the SRS data dictionary.

\textit{How test will be performed:} One team member will read the saved data using a development build or Unity editor. A second team member will compare the fields against the \srstableref{tab:JourneyDataDictionary}{SRS data dictionary} and record any extra fields.

\item[\testdef{ST_Immunity}]\textbf{Dependency Vulnerability Check}

\textit{Type:} Non-functional, Static dependency inspection.

\textit{Requirements Covered:} \srsnfrref{NFR_Immunity}{30}.

\textit{Review procedure:} One team member will list the exact versions of Unity, packages, and plugins included in the release. The member will check these versions against known-vulnerability databases and supplier advisories. A second team member will review any high or critical findings and check that fixes or safe versions have been applied.

\textit{Expected result and evidence:} No included third-party library has a publicly known high or critical vulnerability at the time of release. The team will keep the dependency list, the dated scan or advisory results, and any related GitHub issues.

\end{testlist}


\subsubsection{Cultural and Political Requirements}\label{sec:AMCulturalNFRTests}

The following reviews verify the inclusivity requirement in the SRS
\cite{SRS} and address representation concerns identified in the
Hazard Analysis \cite{HazardAnalysis} through static inspections by
team members rather than usability studies.
Reviewers will examine all released written and visual content,
recording each item's file, screen, or asset ID and build version.


\begin{testlist}

\item[\testdef{ST_AM_InclusivityLanguage}]
\textbf{Inclusivity of Written and Spoken Content}

\textit{Type:} Non-functional, Static content inspection.

\textit{Requirements Covered:} \srsnfrref{NFR_Inclusivity}{31}.

\textit{Review procedure:}

The assigned reviewer and a second team member will separately
review all written and spoken content in the final game,
including instructions, narration, and caregiver guides,
using a checklist of the game's content. For each item, reviewers will answer the
following questions:

\begin{enumerate}[label=\arabic*., leftmargin=3em, labelsep=0.5em, itemsep=0.3em]
    \item Does the content assume that every child has a mother
    and father or a particular family arrangement?

    \item Does the content assign abilities, interests, or
    roles based on gender?

    \item Does the content assume that the user belongs to
    a particular culture, religion, or ethnicity?

    \item Does the content imply that one cultural background
    or family arrangement is normal while others are not?

    \item Does the content contain any stereotypes,
    exclusionary terms, or potentially discriminatory language?
\end{enumerate}

For each question, reviewers will record a Yes/No response.
Any ``Yes'' answer will require a finding containing the content ID,
wording, concern, and proposed change. Reviewers will compare
findings, record their resolution in GitHub Issues, and reinspect
corrected content before release.

\textit{Expected result and evidence:} Every inventory item has a
completed five-question checklist, and there are zero unresolved
inclusivity findings for the reviewed release. The saved checklist,
content inventory, reviewers' names, build version, and GitHub Issue
references provide evidence of the review.


\item[\testdef{ST_AM_InclusivityVisuals}]
\textbf{Inclusivity of Images, Characters, and Rewards}

\textit{Type:} Non-functional, Static visual inspection with manual
in-app confirmation.

\textit{Requirements Covered:} \srsnfrref{NFR_Inclusivity}{31}.

\textit{Initial State:} The final build and complete visual asset
inventory are available; initial setup has been completed.

\textit{Input/Condition:} Inspect every level illustration,
introductory or instructional animation, journey-map image, selectable
character, and cosmetic accessory. Open each available character and
accessory in the app to confirm how it appears in context.

\textit{Output/Result:} No image or character presentation assumes
that all children share one family structure, culture, religion,
gender, or ethnicity. No visual reinforces a stereotype or depicts a
particular identity as the only valid or normal one. All identified
issues are resolved before release.

\textit{Review procedure:}

The assigned reviewer will capture or export the release visuals
and record each asset ID. A second team member will independently
inspect the same visual asset inventory. Both reviewers will answer
the following questions for each asset:

\begin{enumerate}[label=\arabic*., leftmargin=3em, labelsep=0.5em, itemsep=0.3em]
    \item Does the image or character depict one family structure,
    gender, cultural identity, or religious identity as universal?

    \item Does the image or character imply that a person's
    identity determines their abilities, interests, or role?

    \item Are cultural or religious items represented in a
    misleading, inaccurate, or stereotypical manner?

    \item Does the released character library contain
    representations that conflict with the inclusivity
    intent described in the Hazard Analysis?
\end{enumerate}

For each question, reviewers will record a Yes/No response and
explanatory notes. Any ``Yes'' answer will be documented as a
potential concern. The team will track findings in GitHub Issues,
correct the affected assets, and repeat the inspection to
verify that the concerns have been addressed.

\textit{Expected result and evidence:} A complete visual checklist,
screenshots or exported assets, and an issue-resolution log demonstrate
zero unresolved inclusivity concerns at release. The record also
identifies any aspects of diverse character representation that could
not be conclusively assessed by the capstone team.

\end{testlist}


% ============================================================
% 8.8 Compliance Requirements
% ============================================================

\subsubsection{Compliance Requirements}
\label{sec:ComplianceNFRTests}

These tests verify the compliance requirements defined
in the SRS. They check whether the game and its supporting
documentation satisfy the applicable requirements and
standards identified for the project.



\subsection{Traceability Between Test Cases and Requirements}

\wss{Provide a table that shows which test cases are supporting which
  requirements.}

\newpage
\section{Unit Test Description}

This section will be completed once the MIS (detailed design document) has been finalized \cite{MIS}.

% \wss{This section should not be filled in until after the MIS (detailed design
%   document) has been completed.}

% \wss{Reference your MIS (detailed design document) and explain your overall
% philosophy for test case selection.}  

% \wss{To save space and time, it may be an option to provide less detail in this section.  
% For the unit tests you can potentially layout your testing strategy here.  That is, you 
% can explain how tests will be selected for each module.  For instance, your test building 
% approach could be test cases for each access program, including one test for normal behaviour 
% and as many tests as needed for edge cases.  Rather than create the details of the input 
% and output here, you could point to the unit testing code.  For this to work, you code 
% needs to be well-documented, with meaningful names for all of the tests.}

\subsection{Unit Testing Scope}

TBD

% \wss{What modules are outside of the scope.  If there are modules that are
%   developed by someone else, then you would say here if you aren't planning on
%   verifying them.  There may also be modules that are part of your software, but
%   have a lower priority for verification than others.  If this is the case,
%   explain your rationale for the ranking of module importance.}

\subsection{Tests for Functional Requirements}

TBD

% \wss{Most of the verification will be through automated unit testing.  If
%   appropriate specific modules can be verified by a non-testing based
%   technique.  That can also be documented in this section.}

% \subsubsection{Module 1}

% \wss{Include a blurb here to explain why the subsections below cover the module.
%   References to the MIS would be good.  You will want tests from a black box
%   perspective and from a white box perspective.  Explain to the reader how the
%   tests were selected.}

% \begin{enumerate}

% \item{test-id1\\}

% Type: \wss{Functional, Dynamic, Manual, Automatic, Static etc. Most will
%   be automatic}
					
% Initial State: 
					
% Input: 
					
% Output: \wss{The expected result for the given inputs}

% Test Case Derivation: \wss{Justify the expected value given in the Output field}

% How test will be performed: 
					
% \item{test-id2\\}

% Type: \wss{Functional, Dynamic, Manual, Automatic, Static etc. Most will
%   be automatic}
					
% Initial State: 
					
% Input: 
					
% Output: \wss{The expected result for the given inputs}

% Test Case Derivation: \wss{Justify the expected value given in the Output field}

% How test will be performed: 

% \item{...\\}
    
% \end{enumerate}

% \subsubsection{Module 2}

% ...

\subsection{Tests for Nonfunctional Requirements}

TBD

% \wss{If there is a module that needs to be independently assessed for
%   performance, those test cases can go here.  In some projects, planning for
%   nonfunctional tests of units will not be that relevant.}

% \wss{These tests may involve collecting performance data from previously
%   mentioned functional tests.}

% \subsubsection{Module ?}
		
% \begin{enumerate}

% \item{test-id1\\}

% Type: \wss{Functional, Dynamic, Manual, Automatic, Static etc. Most will
%   be automatic}
					
% Initial State: 
					
% Input/Condition: 
					
% Output/Result: 
					
% How test will be performed: 
					
% \item{test-id2\\}

% Type: Functional, Dynamic, Manual, Static etc.
					
% Initial State: 
					
% Input: 
					
% Output: 
					
% How test will be performed: 

% \end{enumerate}

% \subsubsection{Module ?}

% ...

\subsection{Traceability Between Test Cases and Modules}

TBD

% \wss{Provide evidence that all modules have been considered.}

\newpage

\section{Appendix}

\subsection{Symbolic Parameters}\label{sec:VnVParameters}

The first table lists the SRS constants used in these tests. Their values
come directly from the SRS. The second
table defines the few additional values used to run the tests.

\begin{longtable}{@{}p{0.35\textwidth}p{0.12\textwidth}p{0.44\textwidth}@{}}
\caption{SRS Constants}\label{tab:VnVSRSConstants}\\

\toprule
\textbf{Constant} & \textbf{Value} & \textbf{Description}\\
\midrule
\endfirsthead

\multicolumn{3}{c}{\tablename\ \thetable{} -- continued from previous page}\\
\toprule
\textbf{Constant} & \textbf{Value} & \textbf{Description}\\
\midrule
\endhead

\midrule
\multicolumn{3}{r}{Continued on next page}\\
\endfoot

\bottomrule
\endlastfoot
\constdef{ADULT_TEXT_MAX_GRADE} & 8 & Maximum reading grade for caregiver text\\
\constdef{ANSWER_REVEAL_MAX_S} & 1 & Maximum answer reveal time (s)\\
\constdef{AV_SYNC_MAX_MS} & 100 & Maximum audiovisual start offset (ms)\\
\constdef{CLINICIAN_FIND_MAX_MIN} & 1 & Maximum time to find content for a patient (min)\\
\constdef{CLINICIAN_TRAINING_MAX_MIN} & 5 & Maximum clinician training time (min)\\
\constdef{CONTROL_POSITION_TOLERANCE_MM} & 5 & Maximum control position deviation (mm)\\
\constdef{CORE_TASK_MAX_MIN} & 5 & Maximum core task duration (min)\\
\constdef{CPU_UTILIZATION_MAX_PERCENT} & 80\% & Maximum CPU utilization\\
\constdef{CRASH_FREE_RATE_MIN} & 99\% & Minimum crash-free session rate\\
\constdef{FLASH_MAX_PER_S} & 3 & Maximum flashes per second\\
\constdef{FRAME_RATE_COMPLIANCE_MIN} & 95\% & Minimum time meeting frame-rate target\\
\constdef{FRAME_RATE_MIN} & 30 & Minimum frames per second\\
\constdef{HEAT_TEST_SESSION_MIN} & 30 & Minimum device heat-test duration (min)\\
\constdef{INSTALL_SIZE_MAX_MB} & 300 & Maximum installed application size (MB)\\
\constdef{LOW_SPEC_AGE_YEARS} & 5 & Minimum age of low-spec device (years)\\
\constdef{MAINT_FRACTION} & 0.10 & Maximum change effort / initial development effort\\
\constdef{MIN_ANDROID} & 10 & Minimum supported Android version\\
\constdef{MIN_IOS} & 16 & Minimum supported iOS version\\
\constdef{RESPONSE_MAX_MS} & 100 & Maximum time to respond to an input (ms)\\
\constdef{RESPONSE_RATE_MIN} & 95\% & Minimum quick-response rate\\
\constdef{SCREEN_MAX_IN} & 13 & Maximum supported screen diagonal (in)\\
\constdef{SCREEN_MIN_IN} & 4.7 & Minimum supported screen diagonal (in)\\
\constdef{SENTENCE_MAX_WORDS} & 10 & Maximum words per narrated sentence\\
\constdef{SKIP_HOLD_MIN_S} & 1 & Minimum press-and-hold time to skip (s)\\
\constdef{STARTUP_MAX_S} & 5 & Maximum median startup time (s)\\
\constdef{STARTUP_TRIALS_MIN} & 10 & Minimum cold-start trials per device\\
\constdef{TEST_SESSIONS_MIN} & 100 & Minimum test sessions for reliability\\
\constdef{TOUCH_TARGET_MIN} & 15 & Minimum touch target size (mm)\\
\constdef{UNAIDED_RATE_MIN} & 80\% & Minimum unaided task completion rate\\
\end{longtable}

\begin{longtable}{@{}p{0.49\textwidth}p{0.12\textwidth}p{0.30\textwidth}@{}}
\caption{Additional System-Test Parameters}\label{tab:VnVAdditionalParameters}\\

\toprule
\textbf{Constant} & \textbf{Value} & \textbf{Description}\\
\midrule
\endfirsthead

\multicolumn{3}{c}{\tablename\ \thetable{} -- continued from previous page}\\
\toprule
\textbf{Constant} & \textbf{Value} & \textbf{Description}\\
\midrule
\endhead

\midrule
\multicolumn{3}{r}{Continued on next page}\\
\endfoot

\bottomrule
\endlastfoot
\constdef{ANSWER_RETRY_OBSERVATION_S} & 10 & Observation after a reveal with no configured timeout; not an allowed timeout (s)\\
\constdef{TOUCH_TRIALS_MIN} & 40 & Valid touch samples per device\\
\constdef{FREEZE_OBSERVATION_S} & 10 & Time allowed before recording a freeze (s)\\
\constdef{COVERAGE_LINE_TARGET} & 80\% & Line coverage target for game-logic assemblies\\
\constdef{REPLAY_TRIALS_MIN} & 3 & Replays of a completed level when checking for duplicate rewards\\
\constdef{USABILITY_PARTICIPANTS_MIN} & 5 & Participants in the core task usability test\\
\constdef{SURVEY_AGREE_RATE_MIN} & 80\% & Agreement below which a survey statement is reviewed\\
\end{longtable}

\subsection{Reward and Scoreboard Test Data}\label{sec:RewardFixtures}

The Rewards and Scoreboard tests (Section~\ref{sec:RewardsScoreboardFRTests})
use the test data below. The data is loaded only in development builds and
the Unity editor, and is never included in a release build. Each test level
has one skippable exercise and one question with a single correct answer.

\begin{table}[!htbp]
\centering
\caption{Test Levels and Their Rewards}\label{tab:RewardTestLevels}
\begin{tabular}{@{}lrrl@{}}
\toprule
\textbf{Level} & \textbf{Points} & \textbf{Coins} & \textbf{Accessory}\\
\midrule
\constdef{L1} & 10 & 5 & None\\
\constdef{L2} & 20 & 0 & Star badge\\
\constdef{L3} & 40 & 15 & Sun hat\\
\bottomrule
\end{tabular}
\end{table}

\begin{table}[!htbp]
\centering
\caption{Saved States for the Rewards and Scoreboard Tests}\label{tab:RewardSaveStates}
\begin{tabularx}{\textwidth}{@{}lX@{}}
\toprule
\textbf{Saved state} & \textbf{Contents}\\
\midrule
\constdef{SAVE_NEW} & Every level Not Started and 0 coins. No accessories,
  no selected character, and no nickname.\\
\constdef{SAVE_MIXED} & \constref{L1} and \constref{L3} Completed, and
  \constref{L2} In Progress. Coin balance of 3. Owns the Sun hat and the Blue
  scarf, with none equipped. Test character B is selected and the nickname is
  Sam.\\
\bottomrule
\end{tabularx}
\end{table}

\constref{SAVE_MIXED} can be reached by normal play. The child completes
\constref{L1} and \constref{L3} for 20 coins and the Sun hat, buys the Blue
scarf for 17 coins, and then starts and leaves \constref{L2}.

\subsection{Design Verification Checklist}\label{sec:DesignVerificationChecklist}

This checklist will be used during the structured internal review described in
Section~\ref{sec:DesignVerification}. Reviewers will compare the MG and MIS
against the SRS and Hazard Analysis. The checklist is subject to change once the
MG and MIS are finalized.

\begin{itemize}[label={$\square$}, leftmargin=*, itemsep=0.35em]
    \item Are the system's main modules identified in the MG, with a clear
    purpose and responsibility for each one?

    \item Does the MIS provide enough detail about each module's inputs,
    outputs, interfaces, and expected behaviour for another team member to
    implement it?

    \item Are the modules and their dependencies consistent between the MG
    and MIS, without conflicting or unexplained relationships?

    \item Can the functional and non-functional requirements in the SRS be
    traced to the modules or design decisions that support them?

    \item Does the design cover the main gameplay flows, including selecting
    a level, completing it, saving progress, and exiting early?

    \item Does the design account for the safety-related behaviour in the
    Hazard Analysis, such as avoiding punitive feedback, preventing accidental
    input, and stopping audio and animations when a level is exited?

    \item Is the design divided into modules that can be tested and updated
    without unnecessary changes to unrelated parts of the game?

    \item Are external packages and assets identified so their licensing and
    any relevant security or content concerns can be checked?

    \item Are diagrams, interface descriptions, and assumptions clear and
    consistent enough for the team to implement and test the design?
\end{itemize}

\subsection{Usability Survey Questions}\label{sec:UsabilitySurvey}

Participants will complete this survey straight after \testref{ST_Usability} or
\testref{ST_ClinicianUse}, before a team member discusses the session with them.
The survey records how the game seemed to them. It adds evidence for
\srsnfrref{NFR_Layout}{1}--\srsnfrref{NFR_Colour}{10} alongside the measured
results, and does not replace any fit criterion. Participants will answer on
behalf of a child, so the results can identify potential problems but cannot
confirm how children would respond (Section~\ref{sec:Objectives}).

Before the survey, a team member will open the caregiver area so that each
participant can read one text. The form will not record names, email addresses,
or other personal information. Before the first session, one classmate who is
not a test participant will pilot the survey, and any statement they find
unclear will be reworded.

\subsubsection*{Rating Statements}

Participants will rate each statement in Table~\ref{tab:SurveyStatements} from 1
(strongly disagree) to 5 (strongly agree), or choose \textit{Not seen} if they
did not see that part of the game. Participants in \testref{ST_Usability} will
answer as the persona they acted and will skip the clinical staff statements.
Participants in \testref{ST_ClinicianUse} will answer for the parts of the game
they saw and all of the clinical staff statements.

\begin{longtable}{@{}p{0.08\textwidth}p{0.70\textwidth}p{0.14\textwidth}@{}}
\caption{Usability Survey Rating Statements}\label{tab:SurveyStatements}\\

\toprule
\textbf{ID} & \textbf{Statement} & \textbf{Requirement}\\
\midrule
\endfirsthead

\multicolumn{3}{c}{\tablename\ \thetable{} -- continued from previous page}\\
\toprule
\textbf{ID} & \textbf{Statement} & \textbf{Requirement}\\
\midrule
\endhead

\midrule
\multicolumn{3}{r}{Continued on next page}\\
\endfoot

\bottomrule
\endlastfoot
\multicolumn{3}{@{}l}{\textit{Look and Feel}}\\
\surveyq & Controls with the same purpose, such as the exit and hint buttons,
  stayed in the same place on every screen. & \srsnfrref{NFR_Layout}{1}\\
\surveyq & I knew where to look for a control without searching the screen. &
  \srsnfrref{NFR_Layout}{1}\\
\surveyq & The images and animations would not frighten a child aged 3 to 7. &
  \srsnfrref{NFR_ChildSafe}{2}\\
\surveyq & The narration and sounds would not frighten a child aged 3 to 7. &
  \srsnfrref{NFR_ChildSafe}{2}\\
\surveyq & The game showed the procedure in a calm and reassuring way. &
  \srsnfrref{NFR_ChildSafe}{2}\\
\addlinespace[4pt]
\multicolumn{3}{@{}l}{\textit{Usability and Humanity}}\\
\surveyq & The buttons and other touch areas were large enough for a young child
  to touch accurately. & \srsnfrref{NFR_TouchTarget}{3}\\
\surveyq & It was clear how to leave a level at any point. &
  \srsnfrref{NFR_Exit}{4}\\
\surveyq & I could start the level I was asked to find without instructions. &
  \srsnfrref{NFR_Usability}{5}\\
\surveyq & In each exercise, I knew what to do next without help. &
  \srsnfrref{NFR_Usability}{5}\\
\surveyq & The game did not pressure me to keep playing or to come back. &
  \srsnfrref{NFR_NoPressure}{7}\\
\surveyq & The narrated sentences were short enough for a child aged 3 to 7 to
  follow. & \srsnfrref{NFR_Language}{8}\\
\surveyq & The narration used words that a child aged 3 to 7 would know. &
  \srsnfrref{NFR_Language}{8}\\
\surveyq & The text in the caregiver area was easy to understand. &
  \srsnfrref{NFR_Language}{8}\\
\surveyq & The text and icons were easy to see on this screen. &
  \srsnfrref{NFR_Accessibility}{9}\\
\surveyq & I could follow each exercise from the visual cues without relying on
  the narration. & \srsnfrref{NFR_Accessibility}{9}\\
\surveyq & I could tell the level statuses on the journey map apart without
  relying on colour. & \srsnfrref{NFR_Colour}{10}\\
\surveyq & I could tell which answer was correct after it was revealed without
  relying on colour. & \srsnfrref{NFR_Colour}{10}\\
\addlinespace[4pt]
\multicolumn{3}{@{}l}{\textit{Clinical staff only}}\\
\surveyq & Finding a named level would be fast enough during an appointment. &
  \srsnfrref{NFR_ClinicianUse}{6}\\
\surveyq & The training I received was enough for me to show the game to a
  patient. & \srsnfrref{NFR_ClinicianUse}{6}\\
\end{longtable}

\subsubsection*{Open-Ended Questions}

\begin{enumerate}[label=OQ\arabic*:, leftmargin=*, itemsep=0.35em]
    \item Which screen or control was the hardest to find or use, and why?

    \item Which image, sound, word, or sentence might upset or confuse a child
    aged 3 to 7, and why?

    \item Did anything make you feel rushed or pressured to keep playing? If so,
    what was it?

    \item What single change would most help a child use the game without help?

    \item (Clinical staff only) What would stop you from using the game during
    an appointment?
\end{enumerate}

\subsubsection*{Use of the Results}

The VnV Report will present the classmate and clinical staff results separately.
For each statement, it will give the median rating and the number and percentage
of participants who rated it 4 or 5, excluding \textit{Not seen}. A diverging
stacked bar chart will show the responses to each statement, grouped by
requirement. With few participants, these values describe the group and will not
be treated as statistically significant.

The survey does not pass or fail a requirement. If fewer than
\constref{SURVEY_AGREE_RATE_MIN}\% of participants rate a statement 4 or 5, the
team will open a GitHub Issue and review the statement against the related test
results and session notes. Open-ended answers will be grouped by screen, and any
problem reported by more than one participant will also be recorded in GitHub
Issues.

\newpage{}

\section{Appendix --- Generative AI Use Disclosure}
\subsection{AI Tools Used}

\begin{itemize}
    \item ChatGPT (GPT-6.1 Sol, GPT-6)
    \item Claude (Opus 5.5)
\end{itemize}

\subsection{How and Where Used}

\begin{itemize}
  \item In Section~\ref{sec:section2} (General Information), GPT-6 was used to assist
  with organizing, wording, and formatting the Summary, Objectives,
  Custom Documents, and Relevant Documentation subsections.
  The content was reviewed to ensure consistency with the SRS,
  project scope, and decisions made by the team.

  \item In Section~\ref{sec:vnvTeam} (Verification and Validation Team), GPT-6
  was used to improve the organization, wording, and formatting
  of the team responsibilities and testing assignments. It also
  assisted with the LaTeX formatting of the team roles table.
  The assigned roles and responsibilities were based on the
  team's agreed division of work.

  \item In Section~\ref{sec:DesignVerification}, the team developed the design
  verification approach. GPT-6.1 Sol was used to improve the wording, grammar,
  and flow while preserving the team's original ideas and decisions.

  \item In Section~\ref{sec:AMSetupFRTests}, \ref{sec:AMJourneyFRTests}, \ref{sec:AMLevelContentFRTests} (Tests for Functional Requirements), GPT-6
  was used to assist with drafting, organizing, and formatting
  the functional system tests for Setup, Journey Map, and
  Level Content (FR1--FR11). This included structuring the test
  conditions, initial states, expected results, test case
  derivations, and execution procedures. The tests were reviewed
  against the SRS to ensure appropriate requirement coverage,
  traceability, and feasibility within the capstone scope.

  \item In Section~\ref{sec:AMCulturalNFRTests}, GPT-6 was used to assist
  with drafting and formatting the Cultural and Political
  Requirements tests covering NFR31 (Inclusivity). This included
  organizing the static review procedures, inspection checklists,
  expected results, and evidence to be collected. The tests were
  reviewed against the SRS and Hazard Analysis.

  \item In Sections~\ref{sec:PerformanceNFRTests}--\ref{sec:SecurityNFRTests},
  GPT-6.1 Sol was used to draft the non-functional system tests. The team then
  reviewed and revised the test procedures, expected results, and scope remarks
  to align with the SRS and the resources available for the capstone.

  \item In Section~\ref{sec:VnVParameters}, GPT-6.1 Sol was used to help
  organize the symbolic parameter tables and link constants to the tests.
  The team verified the values against the SRS.

  \item In Section~\ref{sec:DesignVerificationChecklist}, GPT-6.1 Sol was used
  to draft the design verification checklist, which the team then reviewed
  and revised.

  \item In Section~\ref{sec:AutomatedTools}, Claude (Opus 5.5) was used to draft
  the tool summary and the coverage plan from the Development Plan and the
  team's issue. The team checked the tool choices against the Development Plan
  and confirmed which tools are still under evaluation.

  \item In Section~\ref{sec:RewardsScoreboardFRTests} (Rewards and Scoreboard),
  Claude (Opus 5.5) was used to draft the functional system tests for
  FR28--FR30 and the test data in Appendix~\ref{sec:RewardFixtures}. It also
  added the \constref{REPLAY_TRIALS_MIN} parameter. The team member directed
  the revisions, such as one test per requirement, automated checks of the
  saved game state, and simpler test data. The tests and expected values were
  checked against the SRS.

  \item In Sections~\ref{sec:LookAndFeelNFRTests} and
  \ref{sec:UsabilityHumanityNFRTests} (Look and Feel, and Usability and
  Humanity), Claude (Opus 5.5) was used to draft the non-functional system
  tests for NFR1--NFR10, following the format of the other non-functional
  tests. It also added the \constref{USABILITY_PARTICIPANTS_MIN} parameter.
  The team member reviewed the procedures and remarks against the SRS fit
  criteria and the capstone limits on testing with children and clinical
  staff.

  \item In Appendix~\ref{sec:UsabilitySurvey} (Usability Survey Questions),
  Claude (Opus 5.5) was used to draft the survey statements, the open-ended
  questions, and the plan for using the results, based on the course checklist
  and NFR1--NFR10. It also added the \constref{SURVEY_AGREE_RATE_MIN} parameter.
  The team member reviewed the survey against the SRS.

  \item GPT-6.1 Sol \& GPT-6 were also used to assist with automatic test numbering,
  cross document references, and general \LaTeX{} formatting.
\end{itemize}

\subsection{Verification of AI-Generated Content}

The team reviewed the AI-assisted content against the SRS, Hazard Analysis,
design documents, and the scope of the project. The proposed verification and
validation activities were checked to ensure that the selected methods were
appropriate, the roles and review procedures were clear, and the planned
evidence would be sufficient to assess the requirements. The system and unit
test plans were checked for requirement coverage, traceability, clearly defined
inputs and expected results, and feasibility within the time and equipment
available to the team. Where a test could not fully meet the SRS fit criterion
within the capstone scope, the limitation was stated explicitly.

The team also checked AI-suggested thresholds, symbolic parameters, and
requirement references against the source documents. Checklists and review
procedures were examined for clarity and completeness. Finally, the \LaTeX{}
document was reviewed for correct numbering, cross-references, tables,
citations, and formatting. The team remains responsible for the accuracy,
completeness, consistency, and appropriateness of the final VnV Plan.

\newpage{}

\section{Appendix --- Reflection}

Questions 1 and 2 were answered individually by each team member. For
Questions 3 and 4, the team brainstormed the knowledge and skills needed to
complete the VnV activities and possible ways to learn them. Each member then
selected one skill and reflected on the approaches they would use to develop it.
The individual responses are provided below.

\subsubsection*{Mohammad Bilal}

\begin{enumerate}
    \item \textbf{What went well while writing this deliverable?}

    One thing that went well was coming up with the overall approach for the
    system tests. When writing the SRS, we tried to make each requirement distinct
    and avoid unnecessary overlap. This was especially helpful for the NFRs,
    which I worked on in both the SRS and this document. Since the requirements
    were already fairly clear and separate, I could plan one test for each NFR
    in my assigned sections. That gave me a good starting point, and I could
    then focus on working out the test conditions, expected results, and how
    each test would actually be carried out.

    Another thing that went well was dividing the work early, as we have done
    for previous deliverables. Unlike the SRS, most sections of the VnV Plan
    did not depend heavily on one another, so team members could make progress
    independently and at their own pace. Assigning responsibilities early also
    made it easier to stay organized and manage the limited time we had for
    the deliverable.

    \item \textbf{What pain points did you experience during this
    deliverable, and how did you resolve them?}

    One challenge was figuring out how we could validate the game with its
    intended users. Direct testing with paediatric patients would require
    Research Ethics Board (REB) approval, and our supervisor explained that
    obtaining this approval would likely be infeasible within the capstone
    timeline. We discussed the issue with our TA and course instructor and
    adjusted our validation plans accordingly. Instead of relying on direct
    testing with children, we will need to use feedback from ophthalmologists
    and other people who work with the intended users, along with simulated
    activities where appropriate. This is not quite the same as testing with
    the children themselves, so it is an important limitation of our plan.

    Another challenge was the number of tests we needed to plan. The SRS has
    43 functional requirements and 29 non-functional requirements, and writing
    a meaningful test for each takes time. It was not enough to simply restate
    the requirement; we had to decide what conditions to test, how to measure
    the result, and whether the test was realistic with the tools and devices
    available to us. AI helped generate some starting ideas, but many of its
    suggestions were too generic or involved testing resources we would not
    have. I had to revise those suggestions and work through the details
    manually to make the tests specific and achievable for our project.

    \item \textbf{What knowledge and/or skill will you need to acquire to successfully complete
    the verification and validation of your project?}

    The team identified several skills that we will need for verification and
    validation, but the one I want to focus on is learning how to use the Unity
    Profiler. Many of the non-functional tests I worked on require actual
    measurements, particularly for frame rate, response time, CPU usage, and
    memory use. I need to learn which measurements to collect, how to record
    them consistently during gameplay, and how to interpret the results. For
    example, I should be able to tell whether a drop in frame rate is a brief
    issue while loading a level or a problem that continues during gameplay.
    I also want to understand the difference between profiling in the Unity
    Editor and profiling a build on a device, since the Editor itself can
    affect performance. Learning this will help me carry out the NFR tests
    properly and provide evidence for whether the game meets the SRS targets,
    especially with the limited physical devices available to our team.

    \item \textbf{What approaches will you pursue to acquire this
    knowledge or master this skill, and why?}

    My first approach will be to work through the official Unity Profiler
    documentation and its examples, focusing on the CPU Usage, Memory, and
    frame timing tools. The documentation should help me understand how to
    capture a profiling session and what the different measurements mean.
    I will then practice using the Profiler on our game as development
    progresses, such as by recording a playthrough of a level and looking at
    what happens when animations or scene transitions occur. This will give
    me experience collecting and comparing results before we run the final
    NFR tests.

    My second approach will be to use AI when I have a specific question or
    run into a problem with the Profiler. For example, it could help me
    understand an unfamiliar measurement, troubleshoot a setup issue, or
    suggest what to check when performance results seem inconsistent. This
    would be a quick way to get unstuck while learning the tool. However, I
    would still check any suggestions against Unity's documentation and test
    them myself rather than assuming the AI explanation is correct. Using
    both approaches should help me learn the Profiler more efficiently while
    making sure the measurements we use in the VnV Plan are reliable.

\end{enumerate}


\subsubsection*{Ahsan Muzammil}

\begin{enumerate}

    \item \textbf{What went well while writing this deliverable?}

    One thing that went well was how quickly our team was able to divide
    the work and establish clear responsibilities for the VnV Plan.
    Early in the process, we identified the sections and requirement
    categories that needed to be completed and assigned them to team
    members through GitHub issues. This made it easier for everyone to
    understand what they were responsible for and what was expected
    of them for the deliverable.

    Using GitHub issues also helped us track everyone's progress and
    monitor how much work each person had been assigned. This was
    particularly useful for ensuring that responsibilities were
    distributed fairly and that no individual was taking on a
    disproportionate amount of work. It also gave us a way to
    identify incomplete tasks and follow up on them before the
    deadline.

    Personally, I was able to complete my assigned sections and
    functional system tests in a timely manner with minimal
    conflicts. Since the responsibilities were clearly defined,
    I could work independently without frequently needing
    clarification from other team members. Overall, our approach
    to task allocation and progress tracking made the writing
    process more organized and efficient.


    \item \textbf{What pain points did you experience during this
    deliverable, and how did you resolve them?}

    One of the main challenges we experienced was a disconnect
    between our understanding of the capstone course expectations
    and our supervisor's expectations regarding communication,
    stakeholder involvement, and the scope of the project.
    This created uncertainty about how we could proceed with
    certain development and validation activities, particularly
    those involving our intended users.

    A significant concern was that directly involving children
    and caregivers in usability testing or interviews would
    require Research Ethics Board (REB) approval. Our supervisor
    explained that pursuing this approval would not be feasible
    within the capstone timeline. As a result, we could no
    longer proceed with the Client Interview Report that we
    had originally considered for our Fall Custom Document.
    We also became uncertain about which Custom Documents
    would be appropriate for the Winter term without access
    to our primary stakeholders.

    To address these issues, we discussed our concerns with
    our TA, Rashad, who recommended that we communicate our
    situation and questions to the course staff for further
    guidance. We also began reconsidering our planned Custom
    Documents, identifying a Cognitive Walkthrough Report or
    Wireframe Report as possible alternatives for the Fall
    semester. For the VnV Plan, we adjusted our proposed
    validation activities to focus on internal reviews,
    peer evaluations using child and caregiver personas,
    and feedback from our supervisor where feasible.

    Although these changes allowed us to continue planning
    the deliverable, some uncertainty remains regarding
    the final selection of our Custom Documents and the
    validation activities we can realistically perform.
    This experience highlighted the importance of
    establishing clear expectations with supervisors
    early in the project and communicating potential
    constraints before they affect planned deliverables.


    \item \textbf{What knowledge and/or skill will you need to acquire
    to successfully complete the verification and validation of
    your project?}

    One skill I will need to develop is the ability to implement
    and execute reproducible functional tests using Unity and
    its testing tools. For the VnV Plan, I was responsible for
    planning the functional system tests covering Setup,
    Journey Map, and Level Content (FR1--FR11), as well as
    the inclusivity requirement (NFR31). While I was able
    to define the test conditions, expected results, and
    execution procedures, carrying out these tests will
    require a better understanding of how to verify
    application behaviour consistently.

    In particular, I want to learn how to use the Unity
    Test Framework and Play Mode tests to verify interactions
    such as starting a preparation level, updating its
    completion status, replaying a completed level, and
    preserving progress between application launches.
    I also need to understand how to create repeatable
    initial conditions, such as resetting saved data
    before a test, so that results are not affected by
    previous gameplay sessions.

    Some of my assigned requirements involve educational
    content and will still require manual inspection.
    Therefore, I also want to improve my ability to
    combine automated functional checks with structured
    manual testing and record evidence such as screenshots,
    test results, and observed behaviour. Developing these
    skills will help ensure that our tests provide reliable
    evidence that \progname{} meets its functional requirements.


    \item \textbf{What approaches will you pursue to acquire this
    knowledge or master this skill, and why?}

    My first approach will be to study the official Unity
    Test Framework documentation and work through examples
    of Edit Mode and Play Mode testing. I will focus on
    understanding how tests are structured, how assertions
    are used to compare actual and expected results, and
    how Unity handles scene loading and application state
    during testing. This will provide a foundation for
    determining which of my assigned functional requirements
    can be verified automatically and which require manual
    interaction on a physical device.

    My second approach will be to practice writing and
    executing tests directly within our \progname{} project
    as development progresses. I plan to begin with
    relatively simple behaviours, such as verifying
    that a preparation level can be selected and started
    regardless of the completion status of other levels.
    I can then work toward testing more complex behaviours,
    including saving progress, preserving completed
    statuses during replay, and resetting application
    data to reproduce first-launch conditions.
    Running these tests against the actual application
    will help me identify limitations in my understanding
    and improve the reliability of the test procedures.

    I will also review my testing approach and results
    with other team members to confirm that the expected
    behaviour matches the SRS and that the tests are
    sufficiently clear for someone else to reproduce.
    Combining documentation-based learning with hands-on
    practice and peer feedback will help me develop the
    skills needed to execute my assigned VnV tests
    effectively and document trustworthy results.

\end{enumerate}


\subsubsection*{Jake Finlay}

\begin{enumerate}

    \item \textbf{What went well while writing this deliverable?}

    Once again, what went best was splitting up the work early using GitHub
    issues. I have said this in previous reflections, but it makes a real
    difference. Everyone knew what was going on and who was doing what, so the
    project moved much more smoothly.

    Another thing that went well was the improved test case template that Ahsan
    committed. It gave me a clear format to follow, which made it much easier
    to write my tests and contribute to both the functional and non-functional
    sections.

    \item \textbf{What pain points did you experience during this
    deliverable, and how did you resolve them?}

    My main pain point was verifying everything against the SRS. Even with AI
    to help us, every test still had to be checked to confirm that it was
    testing the correct requirement. To make this easier, I linked the tests
    directly to the SRS using the \texttt{\#} syntax in our links. Clicking a
    link opens the SRS at the exact position of the requirement. This made it
    much easier to verify what each test covered and which requirements were
    being verified.

    \item \textbf{What knowledge and/or skill will you need to acquire to
    successfully complete the verification and validation of your project?}

    The team identified several skills that we will need, but the one I want to
    focus on is accessibility auditing. One of the tests I planned is the WCAG
    2.2 Level AA audit (\testref{ST_Accessibility}). It requires us to decide
    which success criteria apply to a native mobile game and to justify each
    criterion we leave out. This is not straightforward because WCAG was written
    mainly with web content in mind, and I have not carried out a formal
    accessibility audit before.

    I also need to learn how to check each criterion reliably on our test
    devices. This includes measuring the contrast of text and controls and using
    VoiceOver and TalkBack to confirm that labels are read correctly. I also
    need to understand how these screen readers interact with a Unity interface,
    since Unity does not draw its interface with standard native controls.
    Learning this will help me produce audit results that a second team member
    can reproduce.

    \item \textbf{What approaches will you pursue to acquire this
    knowledge or master this skill, and why?}

    My first approach will be to study the W3C documentation, including the
    Understanding WCAG 2.2 guide and the guidance on applying WCAG to non-web
    software. These are the authoritative sources for what each criterion means.
    They should help me build a list of applicable criteria that I can justify,
    rather than relying on my own interpretation.

    My second approach will be to practise auditing early builds of \progname{}
    as development progresses. I plan to begin with a single screen, using a
    contrast checker and then a screen reader on a physical device. I can then
    compare my findings with those of a second team member and discuss any
    differences. Hands-on practice will show me where the guidelines are hard to
    apply to a game, and it will expose problems while they are still cheap to
    fix. Combining these approaches should help me carry out the final audit
    accurately and give us evidence we can trust.

\end{enumerate}


\newpage
\bibliographystyle{unsrt}

\bibliography{../../refs/References}

\end{document}
